\documentclass[prd, twocolumn, superscriptaddress, nofootinbib,floatfix]{revtex4-2}

\usepackage{graphicx}
\usepackage{amsmath}
\usepackage{amsfonts}
\usepackage{amssymb}
\usepackage{booktabs}
\usepackage[normalem]{ulem}
\usepackage[colorlinks=true, linkcolor=blue, citecolor=blue, urlcolor=blue]{hyperref}

\DeclareRobustCommand{\Eq}[1]{Eq.~\eqref{eq:#1}}
\DeclareRobustCommand{\fig}[1]{Fig.~\ref{fig:#1}}
\DeclareRobustCommand{\figs}[2]{Figs.~\ref{fig:#1} and \ref{fig:#2}}
\DeclareRobustCommand{\app}[1]{App.~\ref{app:#1}}
\DeclareRobustCommand{\sect}[1]{Sec.~\ref{sec:#1}}
\DeclareRobustCommand{\tb}[1]{Table~\ref{tab:#1}}
\DeclareRobustCommand{\tbs}[2]{Tables~\ref{tab:#1} and~\ref{tab:#2}}

\newcommand{\JH}[1]{{\color{blue}JH: {#1}}} %% comments from Jinchen
\newcommand{\SP}[1]{{\color{red}SP: {#1}}} %% comments from Shivam

\graphicspath{{figures/}}

% REVTeX turns on \flushbottom in twocolumn mode, which forces both
% columns to end on the same baseline by stretching the glue between
% paragraphs. On a sparse column that shows up as visibly uneven gaps around
% individual paragraphs. Let the columns end where they end instead.
\raggedbottom

\begin{document}

\title{Diffusion models as preconditioners for Hybrid Monte Carlo in
lattice gauge theory}

%% FILL BEFORE CIRCULATING: each \affiliation below is a placeholder. REVTeX
%% with superscriptaddress attaches an affiliation to every author preceding it,
%% so an author left without one inherits the next author's institution
%% silently. Uncomment and complete all three.
\author{Shivam Panda}
% \affiliation{}
\author{Jinchen He}
% \affiliation{}
\author{Yong Zhao}
% \affiliation{}

\begin{abstract}
Hybrid Monte Carlo (HMC) sampling of lattice gauge fields suffers from critical
slowing down, and its topological charge can freeze entirely at small lattice
spacing. We use a diffusion model as a preconditioner rather than as a sampler:
it performs one inverse renormalization-group step, upsampling a coarse
configuration to a fine one, and HMC is started from that output and remains the
only sampler. The topological charge is transported from the coarse lattice
rather than generated, which is what matters at these couplings, because a
frozen chain keeps whatever charge it is given and cannot correct it within any
practical budget. The network is fully convolutional and conditioned on the
coupling rather than on the lattice size, so a model trained only on lattices up
to $32\times32$ is applied unchanged at $128\times128$. We test the construction in 2D U(1) and U(2). In U(2) it
delivers configurations at $\beta\approx1660$ on a $128\times128$ lattice, where
after $400$ trajectories the preconditioned ensemble agrees with the exact
plaquette to better than $0.2$ parts per million against $25$ for a cold start.
At the benchmark coupling a cold chain under plain HMC does not change
topological sector at all, while every preconditioned ensemble reproduces the
exact charge variance to within $15\%$. Counted in trajectories per independent
configuration, the method is break-even against the strongest classical sampler
these theories admit below $\beta\approx90$, and cheaper by a median factor of
at least nine above it, where that sampler ceases to equilibrate. Widening the
trained coupling range raises the coupling the method reaches, although range
alone does not guarantee it. The cost of building the lifted configuration is accounted
separately and is not a wall-clock advantage in the regime where the classical
baseline is still correct.
\end{abstract}


\maketitle

\section{Introduction}
\label{sec:intro}
Lattice quantum chromodynamics (QCD) provides a nonperturbative formulation of the strong interaction in discretized Euclidean spacetime. In lattice calculations, path integrals are evaluated as ensemble averages over field configurations distributed according to $p(\phi)\propto e^{-S[\phi]}$, where $S$ is the lattice action. These configurations are typically generated using Markov-chain Monte Carlo methods, most commonly Hybrid Monte Carlo (HMC).

Calculations at a finite lattice spacing $a$ are affected by discretization effects, and controlled predictions for continuum QCD require simulations at progressively smaller lattice spacings, followed by a continuum extrapolation toward $a\to0$. This continuum limit is computationally challenging. As the lattice spacing decreases, autocorrelation times grow rapidly because of critical slowing down~\cite{schaefer2011}. The problem is particularly severe for topological observables: the barriers separating topological sectors become increasingly difficult for a continuous Markov chain to cross, leading to topological freezing~\cite{luscher2011}. Over a finite simulation, the chain may then explore only a restricted set of topological sectors and fail to represent the target topological-charge distribution. This also makes initialization increasingly consequential: conventional cold or hot starts of HMC contain no information about the target long-distance or topological structure, so the chain must first thermalize, evolving until its configurations are
distributed according to $e^{-S}$, before they can be used for measurements. When topological tunneling is strongly suppressed, correcting an atypical initial sector population can itself become prohibitively
slow.

In this work, we use a diffusion model as a coarse-to-fine preconditioner for HMC. The model learns to invert a single renormalization-group (RG) step: conditioned on a configuration from a matched coarse lattice, it generates a configuration on the corresponding finer lattice. The coarse configuration provides the long-distance and topological information, while the learned
inverse map restores the short-distance degrees of freedom removed by the RG blocking. The lifted configuration provides an informed initial condition for HMC, which removes the residual mismatch with the fine action and remains the final sampler.

This construction separates the sampling problem across scales. Long-distance fluctuations and topological sectors can be sampled on a coarser lattice where transitions remain accessible, while the diffusion model reconstructs the degrees of freedom required on the finer lattice. The procedure can then be applied recursively, propagating coarse-lattice information toward smaller
lattice spacings where direct HMC becomes increasingly inefficient. We show that this preconditioning reduces the number of trajectories a chain
must run before its configurations are usable from hundreds to tens. As controlled test beds, we study two-dimensional U(1) and U(2) lattice gauge theories. These theories contain compact gauge fields and nontrivial
topological sectors, while U(2) additionally provides a non-Abelian setting.
At the same time, the relevant observables can be compared with exact
analytical results, allowing the complete coarse-to-fine procedure to be
validated directly against the target theory.

Machine-learning approaches to lattice sampling have developed rapidly in recent years. Normalizing flows provide explicit and invertible maps between simple prior distributions and lattice field configurations~\cite{albergo2019,abbott2024practical,abbott2025progress}, while diffusion and score-based models provide an alternative class of generative samplers based on stochastic denoising dynamics~\cite{wang2024,zhu2024,zhu2025,aarts2026}. Related developments have also explored stochastic generative constructions~\cite{wang2026stochastic}. These methods have been extended to non-Abelian and group-valued settings,
through gauge-equivariant architectures and score models defined on the group
manifold~\cite{komijani2026,kanwarvega2025,alharazin2026,vega2025}, to theories
with fermions~\cite{vegaelkhadra2026,tomiya2026fermions}, and to models trained
or extrapolated across couplings and lattice
sizes~\cite{tomiya2026}; see~\cite{wenger2026} for a review of the
four-dimensional SU(3) case.

In most of this related work the learned model is itself part of the sampler:
its output constitutes the generated ensemble, either directly or after an
exactness-restoring correction such as an accept/reject step, as in
Metropolis-adjusted annealed Langevin sampling and sequential Monte Carlo
constructions~\cite{aarts2026,sexty2026}. Generalization beyond the parameters
trained on has also been investigated, in the coupling~\cite{zhu2025,aarts2026}
and in the lattice volume~\cite{tan2026}. In all of these the quality of
sampling at the target parameters depends directly on the accuracy of the model
there. Our use of a diffusion model is different: it constructs only the initial
condition, and the final ensemble is generated by HMC.

A second line of work, more closely related to the present construction, uses RG
information to initialize fine-lattice Markov chains. Coarse-to-fine maps
matched under RG transformations generate fine configurations that are then
rethermalized, reducing the cost relative to cold and hot starts in gauge
theories with and without fermions~\cite{detmoldendres2015,detmoldendres2016},
and the scaling of that advantage toward the continuum has been examined
in~\cite{detmoldendres2018}. Inverse RG transformations have also been used to
drive configurations from small lattices to much larger
ones~\cite{bachtis2022}, and the connection between diffusion models and inverse
RG flow developed in its own right~\cite{cotler2023}. More recently the inverse
RG step has itself been learned, using conditional normalizing
flows~\cite{hasenfratz2026,hasenfratz2026cascade}: the learned map restores the
short-distance modes removed by blocking, a brief rethermalization under the
fine action corrects the remainder, and because the coarse ensemble already
carries the long-distance fluctuations the procedure can be iterated, reaching
large volumes in two-dimensional $\phi^4$ theory. Related hierarchical
constructions include multilevel and super-resolving
flows~\cite{singha2026multilevel,singha2026pss,bauer2025super}.

The present work adopts the same broad separation between coarse sampling, learned upscaling, and fine-lattice rethermalization, but focuses explicitly
on topological-sector structure. In a gauge theory, the coarse configuration contains not only long-distance fluctuations but also information about the
topological sector. Preserving this information during coarse-to-fine generation provides a way to initialize fine-lattice HMC with an appropriate
population of sectors without requiring the fine-lattice chain itself to establish that population through rare tunneling events. This feature is of
particular relevance to the eventual extension to four-dimensional QCD, where topological freezing becomes increasingly severe toward the continuum
limit.

Other machine-learning strategies retain HMC while modifying different parts of the sampling procedure. Learned changes of variables can be inserted inside HMC trajectories to reshape the distribution explored by the integrator, an idea introduced as an analytic trivializing map~\cite{luscher2009} and since
realized with learned transformations~\cite{albandea2023,he2025}, while generative models can also be used to resolve topological sectors separately and determine their relative statistical weights ~\cite{singha2026}. Here we instead leave the HMC evolution unchanged and use the learned model only to construct coarse-informed initial configurations.

The remainder of this paper is organized as follows. In \sect{setup}, we define the lattice theories, notation, and conventional initialization
baselines. \sect{method} introduces the RG-conditioned diffusion preconditioner and the recursive coarse-to-fine procedure.
\sect{seedquality} and \sect{coverage} quantify the quality of the generated initial conditions and examine the training coverage required for
reliable transfer across lattice parameters. \sect{conclusions} then discusses what the construction
would require in four-dimensional lattice QCD.


\section{Simulating gauge fields on a two-dimensional lattice}
\label{sec:setup}

The method is demonstrated in two lattice gauge theories, each chosen because it
freezes topologically in the way four-dimensional QCD does while remaining
exactly solvable, so that every output can be scored against a closed form
rather than against another simulation.
The first is compact U(1), the gauge sector of quantum electrodynamics in two
dimensions. Its compact gauge group gives the field an integer topological
charge and separates the configuration space into sectors. It is also the standard first
testbed for generative models of gauge fields \cite{kanwar2020,zhu2025}. A gauge link carries a group
element on each edge of the lattice; for U(1) that element is a phase,
$U_{x,\mu}=e^{i\theta_{x,\mu}}$ with $\theta_{x,\mu}\in(-\pi,\pi]$, one angle
per edge. We follow the conventions of \cite{gattringerlang2010} throughout.

The second theory is two-dimensional U(2). Its structure,
$\mathrm{U}(2)=(\mathrm{U}(1)\times \mathrm{SU}(2))/\mathbb{Z}_2$, has the
product form of the electroweak gauge group
$\mathrm{SU}(2)_L\times \mathrm{U}(1)_Y$, though the resemblance is structural
rather than physical. The benefit of U(2) lies in that decomposition. As
\sect{freezing} shows, the topological charge is carried entirely by the
determinant, which is an Abelian U(1) factor, while the gauge dynamics remain
genuinely non-Abelian, each plaquette being an ordered product of matrices.
U(2) therefore separates non-Abelian dynamics from non-Abelian topology, and
serves as a computationally simpler test bed on the way to four-dimensional
SU(3).

A U(2) gauge link is a unitary $2\times2$ matrix, which for computational convenience we
store in the split representation
\begin{equation}
    U_{x,\mu} = e^{i\theta_{x,\mu}}\begin{pmatrix}
q_1 + i q_4 & q_3 + i q_2 \\
-q_3 + i q_2 & q_1 - i q_4
\end{pmatrix},
    \label{eq:u2-link}
\end{equation}
an SU(2) element written as a unit quaternion ($\sum_a q_a^2 = 1$) times a
\emph{central} phase, one proportional to the identity and so commuting with
every SU(2) element. Five real scalars therefore describe a U(2) link.

Every observable is built from the plaquette
\begin{equation}
    U_{x,p} \equiv U_{x,\mu} U_{x+\hat{\mu},\nu} U_{x+\hat{\nu},\mu}^\dagger U_{x,\nu}^\dagger,
    \label{eq:plaquette}
\end{equation}
the smallest closed loop, whose trace is therefore gauge invariant; products of adjacent
plaquettes form the larger Wilson loops used as observables. In U(1) the
plaquette reduces to a signed sum of link angles,
\begin{equation}
    \theta_{x,p} \equiv \theta_{x,\mu} + \theta_{x+\hat{\mu},\nu} - \theta_{x+\hat{\nu},\mu} - \theta_{x,\nu},
    \label{eq:theta-plaquette}
\end{equation}
whose order is immaterial; in U(2) the matrix ordering matters. The Wilson
actions,
\begin{align}
    S_{U(1)} &= -\beta \sum_{x}\cos(\theta_{x,p}), \label{eq:action-u1}\\
    S_{U(2)} &= -\beta \sum_{x} \tfrac{1}{2}\mathrm{Re}\,\mathrm{Tr}(U_{x,p}),
    \label{eq:action-u2}
\end{align}
are gauge invariant since they depend only on plaquettes.

\subsection{Topological freezing}
\label{sec:freezing}

The topological charge is the wrapped plaquette sum,
\begin{equation}
    Q = \frac{1}{2\pi}\sum_{x} \mathrm{wrap}\,(\theta_{x,p}),
    \label{eq:q-u1}
\end{equation}
where $\mathrm{wrap}(\alpha)$ returns the representative of $\alpha$ in
$(-\pi,\pi]$.\footnote{Evaluated as the two-argument arctangent of
$(\sin\alpha,\cos\alpha)$, which fixes the quadrant and so the branch.}

The topological susceptibility is the charge variance per unit volume,
\begin{equation}
    \chi_t = \frac{\langle Q^2\rangle}{V},
    \label{eq:chi-t}
\end{equation}
so $\chi_t V = \langle Q^2\rangle$ measures how many sectors an ensemble
populates \cite{bonatirossi2019}.

For U(2), $Q$ is defined by the same formula applied to the determinant plaquette
phase, $\alpha_x = \mathrm{wrap}(\arg\det U_{x,p})$. The central phases in
\Eq{plaquette} are scalars, so they factor out of the product and combine into
exactly the U(1) plaquette sum of \Eq{theta-plaquette}; the SU(2) factors have
unit determinant and drop out of it entirely. Hence
\begin{equation}
    \det(U_{x,p})=e^{2i\theta_{x,p}}.
    \label{eq:det-plaquette}
\end{equation}
The SU(2) component therefore contributes nothing to $Q$. We call the U(1) component of a U(2) link
its determinant sector, $\psi \equiv \arg\det U = \mathrm{wrap}(2\theta)$.

Since $Q$ is defined by the wrapped sum of plaquettes, it can only change if some
$\theta_{x,p}$ is pushed past $\pm\pi$. At large $\beta$ every plaquette angle
clusters near zero due to the Wilson action in \Eq{action-u1}.
Changing $Q$ through purely local updates therefore requires passing through an
intermediate configuration with some $\theta_{x,p}$ near $\pm\pi$, which costs
$\Delta S = 2\beta$ and is accepted with probability $\sim e^{-2\beta}$. That exponential
suppression is the root cause of topological freezing, and it worsens toward the
continuum limit as $\beta$ grows. By \Eq{det-plaquette} the mechanism
is identical in U(2).

\subsection{Classical baseline}
\label{sec:instanton_updates_classical}
In two dimensions, topological freezing can be delayed by a cheap classical
update, and that update is the baseline this paper measures against.

Configurations are generated by HMC \cite{duane1987,neal2011}, discretized by
the Omelyan integrator \cite{omelyan2003} for U(1) and by an exponential-map
leapfrog for the group-valued U(2) fields. The whole trajectory is offered as a
single Metropolis proposal, which is what guarantees sampling from the Boltzmann
distribution despite the integrator's discretization error. That accept--reject
step also makes the chain ergodic at every coupling: given unbounded sampling it
would visit every topological sector and reproduce $P(Q)$ exactly. What fails
toward the continuum is the time this takes. Throughout this paper
\emph{frozen} means that a chain does not traverse sectors within the trajectory
budget it is given, not that those sectors are inaccessible in principle.

\begin{table*}[htbp]
\centering
\small
\caption{The classical baselines, one panel per theory. \emph{Left}, U(1)
steady-state cost at $L=32$ from a cold start, $4$ chains of $3000$
trajectories per arm;
``changes'' counts sector changes in that budget, and ``frozen'' marks an arm
whose charge never moved, for which $\tau_{\rm int}$ is not defined. Plain
periodic HMC never changes sector at any coupling tested, which is why the
freezing diagnostic is a sector-change count rather than $\tau_{\rm int}$.
\emph{Right}, U(2) at $L=16$ from a hot start, $256$ chains and $2000$
trajectories, where flips are changes of the parity of $Q$. The even-only move
cannot change parity, so the flips recorded against it are the ones the
underlying HMC achieves unaided, by pushing a plaquette angle past $\pm\pi$:
frequent at $\beta=14$, rare at $\beta=21$ and absent at $\beta=28$, which is
the local freezing of \sect{freezing} and not a contribution of the winding
update. Its low $\tau_{\rm int}(Q^2)$ is an artifact of never crossing parity at
all, while the marginal move's $2$--$2.7$ is the true cost of sampling that
degree of freedom. Neither panel's $\tau_{\rm int}(Q^2)$ is usable on its own as an ergodicity
test: on the left a frozen arm has none to report, and on the right it reports
healthy on a chain that is not sampling topology.}
\label{tab:classical}
\begin{minipage}[t]{0.52\linewidth}
\centering
\begin{tabular}{r l r c c c}
\toprule
$\beta$ & arm & changes & $\tau_{\rm int}(Q^2)$ & $\langle Q^2\rangle$ & exact \\
\midrule
14.1464 & plain HMC   & 0       & frozen   & 0.0000 & 1.9040 \\
14.1464 & winding HMC & 8{,}447 & 2.85(44) & 1.8705 & 1.9040 \\
55.0237 & plain HMC   & 0       & frozen   & 0.0000 & 0.4743 \\
55.0237 & winding HMC & 5{,}092 & 1.20(13) & 0.4791 & 0.4743 \\
218.58  & plain HMC   & 0       & frozen   & 0.0000 & 0.0290 \\
218.58  & winding HMC & 358     & 1.39(15) & 0.0296 & 0.0290 \\
\bottomrule
\end{tabular}
\end{minipage}\hfill
\begin{minipage}[t]{0.44\linewidth}
\centering
\begin{tabular}{r rc rc}
\toprule
& \multicolumn{2}{c}{even-only} & \multicolumn{2}{c}{marginal} \\
\cmidrule(lr){2-3}\cmidrule(lr){4-5}
$\beta$ & flips & $\tau_{\rm int}(Q^2)$ & flips & $\tau_{\rm int}(Q^2)$ \\
\midrule
14 & 4919 & 0.550 & 72522 & 2.734 \\
21 & 13   & 0.534 & 67298 & 2.365 \\
28 & 0    & 0.548 & 61403 & 1.980 \\
\bottomrule
\end{tabular}
\end{minipage}
\end{table*}

In 2D U(1) a reversible global update that jumps between topological sectors,
accepted by the same Metropolis criterion as the trajectory itself, was
introduced by Albandea \emph{et al.}\ \cite{albandea2021}, who realize it as a
gauge transformation carrying winding around the boundary of a lattice
sub-region. We use the same accept--reject construction with a different
proposal: a field carrying exactly one unit of topological charge, spread as
evenly over the torus as it can be,
\begin{equation}
    I_\nu(x,y) = \frac{2\pi x}{L^2}\,\delta_{\nu,1}
      \;-\; \frac{2\pi y}{L}\,\delta_{\nu,0}\,\delta_{x,L-1},
    \label{eq:instanton}
\end{equation}
the second term being the transition function on the last column, which closes
the field on the torus. Every plaquette angle is then $2\pi/L^2$, so $Q=1$ and
the cost of adding it to a smooth configuration is
$\beta V\big(1-\cos(2\pi/V)\big)\approx2\pi^2\beta/V$. Spreading the flux over
the whole torus is what makes that cost fall with volume, and it is why the
acceptance stays high where local updates are frozen; concentrating the change
on the boundary of a sub-region of size $L_w$ instead costs
$\mathcal{O}(\beta/L_w)$ \cite{albandea2021}. We refer to the move built from
this proposal as winding HMC throughout. It delays the freezing barrier to
much higher $\beta$ than standard HMC without eliminating it
(\tb{classical}).

In 2D U(2) an analogous construction is available because $Q$ depends only on
$\psi$. The SU(2) sector integrates out of the action exactly, plaquette by
plaquette against Haar measure. That integral gives the closed-form marginal weight
$w_{\det}$ of \Eq{wdet}.

The U(1) instanton transfers to U(2) only in part. A purely central phase shift
$\theta\to\theta+\lambda$ commutes with the SU(2) sector entirely, so it can be
applied without disturbing it, at the same cost $\Delta S = 2\pi^2\beta/V$ as in
U(1).
But the determinant sector is $\psi=2\theta$, so the shift moves $\psi$ by
$2\lambda$ and $Q$ by exactly two. Every proposal changes $Q$ by two, so it
cannot change the parity of $Q$, meaning whether $Q$ is even or odd: a chain
under this move stays for its whole history in the parity class it started in.
We call it the even-only move, and a move that crosses between the two classes a
parity flip. Halving the shift does not supply one: it moves $Q$ by one but leaves a
plaquette carrying a spurious $\mathbb{Z}_2$ monodromy, at a cost that does not
fall with volume. The sector structure of two-dimensional $\mathrm{U}(N)$ and
the configurations that carry charge in it are treated in
\cite{rouenhoff2025}. The parity degree of freedom therefore has to be reached
by replacing the SU(2) sector rather than by shifting phases.

Reaching an odd sector from an even one takes three steps: (i) propose
$\psi'=\psi\pm\lambda$, changing $Q$ by $\pm1$; (ii) accept with the Metropolis
ratio of the marginal action alone, $\min(1,e^{-\Delta S_{\det}})$, leaving the
SU(2) sector out of the decision because it is about to be replaced; and (iii)
resample $q\sim p(q\mid\psi')$ exactly by Creutz--Kennedy--Pendleton heatbath
\cite{creutz1980,kennedypendleton1985}. We call this the marginal move.

At $L=64$ and $\beta=416.524$, $200$ cold-started trajectories
over $64$ chains give zero accepted parity flips under the even-only move, so the
measured odd-sector weight is $P(\mathrm{odd})=0$ against an exact $0.4929$.
Under the marginal move the same budget gives $7661$ flips and
$P(\mathrm{odd})=0.4939$. \tb{classical} makes the same comparison at $L=16$ across three
couplings, there from a hot start over $256$ chains.

\section{Method}
\label{sec:method}
Where the baselines of \sect{setup} repair topology inside a running chain, the
method here acts once, before the chain starts. It is one inverse
renormalization-group step, applied repeatedly: a coarse configuration at
$L_c\times L_c$ and coupling $\beta_c$ is lifted to a fine one at $L_f=2L_c$ and
$\beta_f\approx4\beta_c$, with the fine topological charge transported exactly
from the coarse configuration and only the remaining ultraviolet structure
supplied by the model. We call one such $(L,\beta)$ point a \emph{rung}, and the
sequence of rungs running from the base ensemble up to the target the
\emph{ladder}. Before describing the step we fix the direction in which it is
applied, because the topological content of the theory depends on that choice.

The ladder climbs the $(L,\beta)$ plane in one particular direction, and the
choice matters because $P(Q)$ behaves oppositely along the two. Doubling $L$
while quadrupling $\beta$ holds the physical extent $La$ fixed, since
$\beta=1/(e^2a^2)$ sets the spacing \cite{gattringerlang2010}. That is a
continuum limit at fixed physical volume, and along it $\langle Q^2\rangle$
tends to $e^2V_{\rm phys}/4\pi^2$, a finite physical quantity: iterating the tree-level step
$\beta_f=4\beta_c$ out of the deployed U(2) base at $L=16$, $\beta=28$, the
exact determinant-sector value runs from $1.00119$ to $0.92662$ over four such
steps, converging rather than vanishing. Along the deployed ladder, whose
couplings solve \Eq{topology-matching} rather than the tree-level rule, it is
instead constant to six digits. Transporting the charge is therefore the correct operation along the
ladder, and $P(Q)$ is the same distribution at every rung up to lattice-spacing
corrections. Raising $\beta$ at fixed $L$ instead shrinks the physical volume and collapses
the charge distribution with it, so a charge carried along that direction
unchanged would be wrong: at $L=32$ in 2D U(1) the exact
$\langle Q^2\rangle$ falls from $7.74$ to $0.014$ as $\beta$ runs from $4$ to
$256$. What that direction removes is the volume, not the topology. The coarse
ensemble's own distance from the continuum remains, and it is the one systematic
the construction inherits rather than creates: a finer base would carry a
smaller one. \sect{coverage} evaluates the fixed-$L$ direction instead, one
independent lift per coupling, and never carries a charge across it.

\subsection{Training data}
\label{sec:naive-blocking}


Blocking maps a fine configuration to a coarse one at negligible cost, but it
is not invertible: many fine
configurations block to the same coarse one, and the task here is to draw from
that fibre at the correct weight. That is what the model supplies, and it is
also why a blocking map alone would not do. HMC repairs short-distance
structure quickly, so a crude inverse would be enough if local observables were
all that mattered; what HMC cannot repair within any practical budget is the
topological sector, and the sector is exactly what the coarse configuration
already carries.

A fine configuration is generated at a $\beta$ small enough that winding HMC has
a finite decorrelation time even for the slow modes. Its coarse partner is fixed completely by summing the neighboring links,
\begin{align*}
\theta_x'(X,Y) &= \mathrm{wrap}\big(\theta_x(2X,2Y) + \theta_x(2X+1,2Y)\big),\\
\theta_y'(X,Y) &= \mathrm{wrap}\big(\theta_y(2X,2Y) + \theta_y(2X,2Y+1)\big),
\end{align*}
with $\theta'$ the coarse link angle. In U(2) the same map acts on $\psi$, and
the statement survives the non-Abelian group for one reason: the determinant is
a homomorphism, so the determinant phase of a product of links is the plain sum
of their phases, and \Eq{det-plaquette} telescopes exactly as the Abelian case
does. This makes the coarse plaquette angle equal
the wrapped sum of its four constituent fine sub-plaquettes exactly, and hence
$Q_{\mathrm{coarse}} = Q_{\mathrm{fine}}$.
Nothing downstream depends on that equality holding exactly: at deployment the
fine charge is imposed on the lifted configuration rather than inherited from
the blocking map.

The blocked coupling must be fixed by requiring the two lattices to agree on
some observable, as in the Monte Carlo renormalization group
\cite{swendsen1979}, and here the choice is not free. Sector
transport is an identity, so whatever $P(Q)$ the base ensemble carries is what
every rung above it receives, and the number of sectors it populates is
$\chi_t V$ of \Eq{chi-t}. We therefore fix $\beta_f$ at each step as the root
of
\begin{equation}
    \langle Q^2\rangle_{\beta_f,\,L_f} \;=\; \langle Q^2\rangle_{\beta_c,\,L_c},
    \qquad L_f = 2L_c,
    \label{eq:topology-matching}
\end{equation}
with both sides evaluated in the closed form of \app{exact}. A solution exists
because to leading order $\chi_t V\approx V/(4\pi^2\beta)$, so doubling $L$
quadruples $V$ while quadrupling $\beta$ restores it, and the root sits near
that tree-level estimate; \app{matching} gives the numerical solution and the
resulting schedules.

Since every rung inherits the base's $P(Q)$, the base must sample it correctly,
and it is cheap enough to check. The deployed U(2) base is $4096$ configurations
over $1024$ chains at $L=16$, $\beta=28$, generated by HMC with the marginal
winding move at an acceptance of $0.35$ and with no sector seeding of any kind.
Against the closed form it gives $\langle Q^2\rangle=0.967$ against $1.00119$
and $P(\mathrm{odd})=0.512$ against $0.493$, with the sector weights tracking
the exact ones from $Q=-3$ to $+3$; neither deviation is resolved at this
statistics. The even-only move records no parity flips at this coupling
(\tb{classical}), which is why the base is built with the marginal one.

At high $\beta$ the fine ensemble used to build training pairs is itself frozen
in $Q$, with $\langle Q^2\rangle\ll1$, so HMC returns configurations almost
entirely at $Q=0$. A training set built purely from HMC output at a fixed high
coupling would then never show the network the local structure of a
configuration at nonzero $Q$ at that coupling.
We manufacture the missing coverage rather than sampling it. Sector
augmentation takes a fixed fraction of the fine ensemble at a rung, draws a
charge shift uniformly from a small integer range, adds that many copies of the
instanton field of \Eq{instanton}, and relaxes the result
locally at fixed topology to remove the ultraviolet error the instanton leaves
behind. It consults no closed form, so nothing in the construction of the
training data asks the theory to be solvable.

\subsection{Score network}
\label{sec:scoring-network}

\begin{figure*}[bt]
  \centering
  \includegraphics[width=0.9\linewidth]{u1_2d/45_architecture.png}
    \caption{The gauge-covariant score network, shown for U(1), with
  $324{,}954$ parameters in $4$ FiLM blocks of width $56$. The U(2) network has
  the same structure, with width $64$ and $403{,}394$ parameters. The network outputs
  one number per plaquette and assembles the score from them with a lattice
  curl; \sect{scoring-network} shows that this costs no generality, and that
  nothing in the architecture refers to the lattice size.}
  \label{fig:architecture}
\end{figure*}

Diffusion models generate by reversing a noising process. A forward It\^o SDE
$dx = f(x,t)dt + g(t)dW$ carries the data distribution to a tractable reference
over $t\in[0,T]$. By Anderson's theorem \cite{anderson1982} it has a
reverse-time companion with the same marginals, whose only unknown is the score
$\nabla_x\log p_t(x)$; denoising score matching \cite{vincent2011} estimates
that score by regressing on the known conditional transition kernel, and Song et
al.\ \cite{song2021} supply the discretizations used here. The standard
construction assumes $x\in\mathbb R^n$ with a Gaussian kernel. Neither holds for
link angles, which live on a compact periodic domain, so both the forward
process and the score it induces are built from the domain's own kernel.

Each configuration is a vector $\vec\theta$ with $2L^2$ components, noised
independently with the wrapped Gaussian transition kernel
\begin{equation}
    K(\theta;\hat\theta,\sigma_t)=\frac{1}{\sigma_t\sqrt{2\pi}}
      \sum_{r=-\infty}^{\infty}
      \exp\!\left(-\frac{(\theta-\hat{\theta} + 2\pi r)^2}{2\sigma_t^2}\right),
    \label{eq:wrapped-kernel}
\end{equation}
with $\hat\theta$ the clean angle. As $\sigma_t$ grows this approaches the uniform
distribution on the circle, and the regression target is $\nabla_\theta\log K$.

The architecture is shown in \fig{architecture}. The network sees two
fields, the noisy fine configuration and the coarse configuration conditioning
it, each as gauge-invariant channels and never as raw link angles. The fine
field contributes six: the cosine and sine of its plaquette angle and of the two
$1\times2$ rectangle angles, which are the smallest loops that distinguish the
two lattice directions. The coarse field contributes five: the cosine and sine
of its plaquette angle and of its $2\times2$ loop angle, together with the
wrapped plaquette angle itself, which a cosine and sine pair cannot distinguish
from its negative and which is what carries the sign of the local charge
density. Eleven channels enter the first convolution. Because the two fields
live on grids of different size, the coarse one is first enlarged to the fine
grid by nearest-neighbor replication, each coarse value copied onto all four
sites of its block, after which the two are concatenated channel-wise. The noise level $\sigma_t$ and
coupling $\beta$ enter through a sinusoidal embedding feeding per-block
Feature-wise Linear Modulation (FiLM) \cite{perez2018}, into which the
spatial mean of the coarse conditioning channels is also projected. That gives
every block direct access to a genuinely global quantity, since the mean wrapped
coarse plaquette angle equals $2\pi Q_{\mathrm{coarse}}/V$ by definition
of $Q$, and a convolutional network with bounded kernel size cannot otherwise
reconstruct a lattice-wide sum.

Gauge invariance is built into the architecture, by the construction of
\Eq{curl-x} and \Eq{curl-y} given
below. The Wilson action's $90^\circ$-rotation, reflection and
charge-conjugation symmetries are not. Training instead applies a uniformly
random element of that $16$-element group, before blocking and noising, to a
randomly selected half of the mini-batches, one element shared across the
batch.

The network predicts one scalar $h_p$ per plaquette rather than a link-valued
score, and the score is assembled by the lattice curl,
\begin{align}
  s_x(x,y) &= h_p(x,y) - h_p(x,y-1), \label{eq:curl-x}\\
  s_y(x,y) &= h_p(x-1,y) - h_p(x,y). \label{eq:curl-y}
\end{align}
This is the general form of the gradient of any sum of functions of plaquettes:
every such gradient is a lattice curl of some per-plaquette scalar, and every
per-plaquette scalar produces one. It is therefore complete for the score of the
Wilson density itself, and incomplete in one direction for the conditional the
model actually learns, since blocking fixes the coarse Polyakov loops from the
fine ones while a curl carries no component along those holonomies. No
observable reported here is sensitive to that direction, all of them being
functionals of plaquettes, but it is a limitation of the head rather than of the
theory. Each link
belongs to exactly two plaquettes with opposite orientation, so
$\partial\theta_p/\partial\theta_\ell\in\{+1,-1,0\}$ and only two terms survive.
Because $h_p$ is itself a function of gauge-invariant inputs, the score is
gauge-covariant by construction. Nor does any part of the network refer to the
lattice size: the convolutions are circular and the normalization acts on each
site separately, so one trained network applies at every volume on the ladder,
which is what \sect{u1obs} tests directly.

The analytic Wilson force drawn alongside the learned head in
\fig{architecture} is the blended score term of \sect{reverse_diffusion}, gated
on the noise level and used in U(1) only; \app{blend} gives the U(2) measurement
that led to its being switched off there.

Network capacity, the training budget and the couplings each checkpoint was
trained on, together with the sampler and local-sweep settings used
throughout, are collected in \app{settings}.

\subsection{Reverse diffusion process}
\label{sec:reverse_diffusion}

\begin{figure*}[bt]
    \centering
    \includegraphics[width=\linewidth]{reverse_diffusion_schematic.png}
    \caption{The reverse-diffusion chain for 2D U(1) and U(2)
    (\sect{reverse_diffusion}). It runs from right to left, starting from angles
    $\phi_T$ drawn uniformly on the circle and ending at the lifted configuration
    $\phi_0$, with $\phi$ standing for whichever angle the model denoises: the
    link angle $\theta$ in U(1) and the determinant phase $\psi$ in U(2). Here
    $s_\theta$ is the learned score and
    $\Delta_t=\sigma_t^2-\sigma_{t-1}^2$, and each denoising step is followed by
    one Langevin correction step (dashed).
In U(2) the SU(2) sector is then drawn
    conditionally by the exact heatbath of \sect{u2-extension}. The delivered
    configuration additionally passes through a short run of exact local sweeps,
    which is what \sect{reverse_diffusion} separates from the raw lift.}
    \label{fig:method_schematic}
\end{figure*}

The wrapped-Gaussian kernel is exactly the transition kernel of a driftless
diffusion wrapped onto the circle, so $f\equiv0$ and
$g(t)=\sqrt{d\sigma_t^2/dt}$: a variance-exploding schedule in the terminology of
\cite{song2021}. How such a schedule on a group manifold controls the decay of
the action along the diffusion is examined in \cite{komijani2026noise}. The forward process is $d\theta = g(t)dW$ and its reverse-time
companion is
\begin{equation}
    d\theta = -g(t)^2\,s_\theta(\theta_t,\sigma_t,\beta)\,dt + g(t)\,d\bar W,
    \label{eq:reverse-sde}
\end{equation}
so the drift of the reverse process is the learned score. In U(1) one further
term enters: the analytic Wilson score of the target action, evaluated at a
noise-smeared coupling and blended in with a weight that vanishes at large
$\sigma$ and approaches unity at small $\sigma$, so that the network supplies
the drift while the field is noisy and the known action takes over as it becomes
clean. The U(2) sampler uses the learned score alone, and not by omission: the same
blend is implemented there, is deployed at zero weight, and degrades every
observable when it is switched on, for reasons that are not settled
(\app{blend}). Starting from pure noise
on the fine lattice, ancestral sampling integrates this backward through the
noise schedule (\fig{method_schematic}). Every result reported here takes one
Langevin corrector step after each ancestral step \cite{song2021}.

Three mechanisms softly bias $Q_{\mathrm{fine}}$ toward $Q_{\mathrm{coarse}}$
during sampling: the coarse-conditioning channels, the FiLM-embedded global
mean, and a reconstruction-guidance penalty \cite{ho2022} added to the score at each noise
level. That penalty is the gradient, with respect to the cell sum, of
\begin{equation}
    -\frac{\big(\sum_{p=1}^4\theta_{p,\mathrm{fine}}
        - \theta_{\mathrm{coarse}}\big)^2}{2v(\sigma)},
    \qquad v(\sigma)\approx 8\sigma^2,
    \label{eq:guidance}
\end{equation}
in which $v(\sigma)$ is the noise that the eight boundary links of a coarse
cell inject into the blocked plaquette at level $\sigma$. The per-cell scalar is broadcast to the four fine sub-plaquettes and
assembled into link space by the same discrete curl used for the score head, so
the guidance is gauge-covariant by construction. None of these make $Q_{\mathrm{fine}}=Q_{\mathrm{coarse}}$ exact, they only
make the sampler more likely to settle in the correct sector on its own.

The equality $Q_{\mathrm{fine}}=Q_{\mathrm{coarse}}$ is instead imposed
exactly by a fourth step: the sampler computes
$\Delta Q = Q_{\mathrm{coarse}} - Q_{\mathrm{fine}}$ and adds $\Delta Q$ copies
of the instanton field of \Eq{instanton}. It does this at every tenth reverse-SDE step once the noise level
has fallen below $\sigma=0.5$, and once more unconditionally after sampling
completes.\footnote{The threshold is a stated setting rather than a fitted one: above it the
charge of the noisy sample is not yet a meaningful quantity to project onto, and
$0.5$ is where it becomes one. Applying the projection partway through the
reverse process rather than only at the end lets the remaining diffusion steps
and the local sweeps relax the ultraviolet error it leaves behind.}
Sampling is followed by a short tail of exact local updates, one heatbath pass
over the lattice and two overrelaxation passes per sweep, at the counts given in
\app{settings}. Both are exact updates of the fine action at fixed link
topology: each visits one link at a time and draws it from its own conditional
weight, which at these couplings carries a plaquette angle past $\pm\pi$ with
negligible probability, so $Q$ is left as the projection set it. We call this
tail the \emph{local sweeps}, and reserve \emph{rethermalization} for an HMC
chain restarted from a lifted configuration, which is the sense
\sect{relaxation-correction} gives it and the sense of
\cite{detmoldendres2015}. The distinction matters because the two repair different things and only one of
them is a cost the user pays per chain. The resulting convention for what is
scored is stated in \sect{seedquality}.

The charge a preconditioned chain begins with is therefore inherited rather than
sampled, and the argument does not require that inheritance to be exact. HMC is
ergodic (\sect{instanton_updates_classical}), so any distortion of the sector
weights would eventually decay; transport does not replace that guarantee, it
avoids having to spend the trajectories that realizing it would cost. The chain
starts in the sector distribution it would otherwise have had to reach, and the
accept--reject step remains responsible for correctness throughout.

\subsection{Extension to U(2)}
\label{sec:u2-extension}

The construction carries over to U(2) essentially unchanged. The lifting step is a
conditional diffusion model that generates only the determinant sector $\psi$,
with $\psi$ standing in for $\theta$ throughout. Integrating SU(2) out of one
plaquette against Haar measure is the standard one-link integral
\cite{creutz1980}, and gives the exact marginal weight of the determinant-sector
plaquette angle $\alpha$,
\begin{equation}
    w_{\det}(\alpha;\beta) = \frac{2I_1(z)}{z},
    \qquad z=\beta\cos(\alpha/2).
    \label{eq:wdet}
\end{equation}
 This is not of Wilson form, so the raw coupling $\beta$ is not the right thing
to condition the network on. It is conditioned instead on $\beta_1$, the Wilson
coupling whose single-plaquette weight is closest to $w_{\det}$ under the
minimum Kullback--Leibler projection onto the Wilson family
\cite{amarinagaoka2000}. Because the Wilson weight is a von Mises distribution in $\alpha$ with
concentration $\beta_1$, that projection reduces to a moment condition on
$r_1(\beta)=I_1(\beta)/I_0(\beta)$, the normalized first character coefficient
and the mean plaquette at infinite volume,
\begin{equation}
    \frac{I_1(\beta_1)}{I_0(\beta_1)} = \langle\cos\alpha\rangle_{w_{\det}},
    \label{eq:vonmises}
\end{equation}
the standard concentration estimate for that family \cite{mardiajupp2000}, with
$I_\nu$ the modified Bessel function of order $\nu$ and the projection itself
tabulated in \app{exact}. Along the deployed ladder the raw couplings
$\beta=28$, $105.65$, $416.52$ and $1660.08$ project to
$\beta_1=7.02$, $26.42$, $104.13$ and $415.02$. This is a projection between weights on
one lattice, and is unrelated to the ladder matching of
\Eq{topology-matching}, which relates two lattices.

Two couplings therefore appear throughout. The \emph{raw} coupling is
$\beta_{U(2)}$ itself, the coefficient in the U(2) Wilson action of
\Eq{action-u2}; the \emph{model} coupling is $\beta_1$, the matched
U(1) value the score network is conditioned on. Every coupling quoted in this paper, including every training ceiling, is a raw
$\beta$; the model coupling appears only where it is named as such. The
exception is one checkpoint's name, \texttt{cov60}, which was fixed during the
experiments and refers to a model coupling: its rungs stop at $\beta_1\approx57$,
or $\beta=227$ raw.

The SU(2) sector needs no model. It cannot be generated on its own, because the
U(2) action does not factorize into independent U(1) and SU(2) parts, so $q$
must be drawn conditioned on $\psi$. At fixed $\psi$ that conditional is the
standard SU(2) single-link weight,
$p(q \mid \psi) \propto \exp(\beta\, k \cdot q)$. Here $k$ is the U(2) staple
sum, rotated by the link's own phase and written as a quaternion, and that
rotation carries the whole dependence on $\psi$. A fresh SU(2) sector can therefore be equilibrated exactly by
Creutz--Kennedy--Pendleton heatbath sweeps
\cite{creutz1980,kennedypendleton1985}, which leave $\psi$, and hence $Q$,
unchanged.

\section{Preconditioner quality}
\label{sec:seedquality}
What the preconditioner is worth is measured in two ways: how close the
configurations it delivers are to the exact target distribution, and what a chain
started from one costs. We call the output of one inverse
renormalization-group step a \emph{lifted configuration}, and reserve
\emph{preconditioned} for a chain, arm or ensemble started from one. Two objects
are scored, and which one is meant depends on the question. The \emph{raw lift}
is the reverse-diffusion output together with the conditional SU(2) heatbath,
before any local sweep; the \emph{delivered ensemble} is what the local sweeps
of \sect{reverse_diffusion} return. Numbers that
describe the map itself, as in \sect{coverage}, are measured on the raw lift, and
numbers that describe what the pipeline hands to HMC are measured on the
delivered ensemble.
Throughout this section an \emph{arm} is one complete run: a starting
configuration, which is cold, hot or preconditioned, advanced by one named
sampler for a stated trajectory budget. Arms that differ only in their starting
configuration are always advanced by the same sampler for the same trajectory
budget, except in the U(1) volume scan of \sect{u1obs}, whose chain counts
differ as noted there.

Wilson loops from the plaquette up to $8\times8$, probing both ultraviolet (UV)
and infrared structure, together with $P(Q)$ and the topological susceptibility $\chi_t$, all
follow from the single-plaquette character coefficients of each theory.
\tbs{u1-exact-observables}{u2-exact-observables} in \app{exact} collect them. Local observables and topology are
scored separately, each in its own unit. A Wilson loop is quoted as its relative
deviation from exact, $(\text{mean}-\text{exact})/\text{exact}$, in parts per
million (ppm). A charge variance is quoted as a ratio instead,
$\langle Q^2\rangle$ over its exact value, because the arms span the full range
from a frozen chain, where it is zero, to a hot start, where it is $88$ times
exact; a relative deviation in parts per million is unreadable across that
range, while a ratio places the frozen and the correct arm at $0$ and $1$.

\subsection{Two timescales}
\label{sec:relaxation-correction}
Two timescales are needed, both in the form used in the multiscale
equilibration studies of \cite{detmoldendres2015,detmoldendres2016}. The first is the integrated
autocorrelation time $\tau_{\rm int}(\mathcal{O})$, estimated by Sokal's
automatic windowing \cite{madrassokal1988,wolff2004} as specified in
\app{stats}. A run of $N$ measurements carries
$N/2\tau_{\rm int}$ independent samples, so one independent configuration costs
$t_{\rm interval}=2\tau_{\rm int}$, which we maximize over the scored
observables and over $Q^2$ so that it is set by the slowest mode. Including
$Q^2$ is deliberate: an arm whose charge never moves has no finite interval,
which is the correct statement that it supplies no independent configurations of
the target, while for winding HMC $Q^2$ does decorrelate, with
$\tau_{\rm int}(Q^2)$ between $1.2$ and $2.9$ (\tb{classical}). A series that
does not move at all is reported as infinite rather than integrated, since a
constant has no autocorrelation to measure.

The approach to equilibrium is governed instead by the exponential
autocorrelation time $\tau_{\rm exp}$, and how long it takes depends on the
overlap of the starting distribution with that slowest mode
\cite{detmoldendres2015}. Endres \emph{et al.}\ obtain it by fitting that decay
jointly across several starting distributions with a shared exponent, the
natural estimator when the start is far from equilibrium. It is not available
here, because a preconditioned chain can begin at the target: there is then no
transient to fit, and the fit either fails or returns a value set by the window
rather than by the chain. We therefore read the equilibration time off the data
directly.
$t_{\rm therm}$ is the first trajectory after which the bias against the exact
value stays within two standard errors (\app{stats}) for five consecutive
measurements. It is maximized over the Wilson loops but not over
$Q^2$, since a frozen chain has an exactly constant $Q^2$, which would make
every classical equilibration time infinite and the comparison empty; topology
is scored separately, against the closed form. Both limits behave: a start
already at the target returns zero, and one that never arrives returns infinity.
Both timescales are quoted with the chain-resampling interval of \app{stats},
which for $t_{\rm therm}$ carries its dominant uncertainty.

The two timescales enter the cost of an ensemble differently. A classical chain
equilibrates once and then yields one configuration per interval, so its
per-configuration cost tends to the interval as the ensemble grows. A
preconditioned ensemble is built the other way round: each configuration
descends from an independent coarse one, so successive members owe no interval
to one another, and the cost is the equilibration each of them pays. Comparing
the two gives the improvement factor,
\begin{equation}
    F = \frac{t_{\rm interval}\ (\text{classical})}
             {t_{\rm therm}\ (\text{preconditioned})},
    \label{eq:improvement-factor}
\end{equation}
and it is what we report throughout.

\subsection{U(1)}
\label{sec:u1obs}
Two experiments support this subsection. The first compares \emph{entry cost},
what each route pays to reach a usable configuration: the delivered product, one
lift plus the sixteen local sweeps of \app{settings} and no trajectories at all,
against a winding HMC chain through a burn-in of $500$ to $8000$
trajectories. The samplers differ because the entry costs are what is being compared. HMC is
the classical baseline throughout this paper because it is the algorithm of the
four-dimensional target, not because it is the most efficient sampler available
in two dimensions: for local relaxation in a pure gauge theory the heatbath and
overrelaxation sweeps that the delivered pipeline itself uses as a repair tail
are exact and cheaper per unit work than a trajectory. The comparison that
isolates the topological advantage from the local one, a classical start
followed by those same sweeps, is not run here. The second is a volume scan, in which the raw lift, a cold start and a hot
start are all advanced by plain HMC. That scan is matched in sampler but not in
statistics: the preconditioned arm carries $128$, $96$ and $64$ chains at
$L=32$, $64$ and $128$ against the classical arms' $32$, $16$ and $8$. Because
$t_{\rm therm}$ is a crossing in units of a standard error, the better-resolved
arm is held to the tighter band, so the asymmetry works against the
preconditioner and the equilibration advantage quoted below is a lower bound on
what a matched comparison would give.

In the first, the delivered configuration sits inside the exact plaquette's
error band at $t=0$, before any trajectory, at every coupling tested
($\beta_f = 4.44$ to $218.58$, a $50\times$ range). The classical baseline of
record, winding HMC, reaches the same band within a one-time
burn-in of $500$ trajectories at the two easier couplings ($\beta_f=4.44$,
hot-started, and $14.1464$), but not at the two largest. At $\beta_f=55.02$ it
fails the quality criterion after $500$ trajectories and after $2000$, and
passes at $8000$; at $\beta_f=218.58$ it fails at $500$, at $2000$ and at
$8000$ (\tb{u1_h2h}). The classical entry cost therefore rises by at least an
order of magnitude across this range, while the lift's is zero throughout.
These four couplings are an independent $L=16\to32$ grid rather than rungs of
the ladder of \sect{naive-blocking}: each is a single lift, scored on its own,
and no charge is carried between them.

The map itself survives a change of volume. In the volume scan, at fixed
$\beta_f=14.1464$ over $L=32$, $64$ and $128$, the raw lift's plaquette sits
$869$, $892$ and $820$ parts per million from exact before any trajectory or
any local sweep. That flatness is what the architecture predicts: the
network is fully convolutional and conditioned on the coupling rather than on
the lattice size, so nothing in it ties the map to the volumes it trained on.

The sixteen local sweeps of the delivered pipeline remove that deviation, in
keeping with the division of labor set out in \sect{reverse_diffusion}. They
run at fixed topology by construction, so they cannot alter
$\langle Q^2\rangle$: whatever charge the lift delivers is what HMC receives.
Against exact, $\langle Q^2\rangle$ comes out $0.95\pm0.10$ and $1.20\pm0.18$
times exact at $L=32$ and $64$, and $0.71\pm0.13$ at $L=128$. The first two are
consistent with exact and the third is low by $2.3$ standard errors, the only
volume in this scan not reproduced within its error; at $L=128$ the lift draws
on a coarse ensemble four doublings above the base, and we do not resolve
whether the deficit is the coarse ensemble's or the map's. Under the plain HMC of this scan the charge never moves at any of the three
volumes, so that arm's $\langle Q^2\rangle$ vanishes identically. Winding HMC is
the comparison that matters, and it does sample topology here, changing sector
$8447$ times in $3000$ trajectories at this coupling and reaching $0.98$ times
exact (\tb{classical}) against the lift's $0.95\pm0.10$ at $L=32$. On this
observable the two are therefore indistinguishable at this coupling; what
separates them is cost rather than accuracy (\sect{cost}).

The advantage in equilibration cost does not survive equally. Over the same
three volumes, with plain HMC in every arm so that the starting configuration is
the only difference between them, the preconditioned chain needs $9$, $18$ and
$52$ trajectories to equilibrate against a cold start's $168$, $400$ and $528$,
factors of $19$, $22$ and $10$. That the cost grows with volume while the
deviation does not is a property of the criterion rather than of the map:
$t_{\rm therm}$ is a crossing in units of the standard error, which shrinks as
the volume grows, so a deviation of fixed relative size takes longer to fall
inside the band. The cold arm lengthens with volume for the same reason.

\subsection{U(2)}
\label{sec:u2obs}
The U(2) benchmark is where the two failure modes separate most cleanly, because
at $L=64$, $\beta=416.5$ the classical arms fail at UV structure and at topology
independently, and no classical arm repairs both. \tb{u2_observables} gives each
arm's mean Wilson loop after $400$ trajectories on $64$ chains, averaged over
the settled half of every history and compared with the closed form; the exact
$\langle Q^2\rangle$ is $1.0012$. Every arm started from a preconditioned
configuration is correct on both counts, whichever sampler advances it: the
plaquette and $W(2\times2)$ are indistinguishable from exact, deviating by at
most $1.4$ and $10$~ppm against standard errors of $2.2$ and $11$~ppm, while
$W(8\times8)$ is within $1300$~ppm, and $\langle Q^2\rangle$ falls between
$0.98$ and $1.14$ times exact. The upper end is
the plain-HMC arm, whose charge is inherited from the preconditioner and cannot
move, so that $1.14$ is a single frozen value rather than a sampled distribution. No
classical arm reaches both the UV and the topological target.

The uncertainties change how the two volumes should be read. At $L=64$ every
lifted arm is resolved against exact on the plaquette and $W(2\times2)$ and
agrees there; at $L=32$ the lattice is a quarter the size and the large loops
carry a per-configuration spread wide enough that none of the lifted arms'
$W(8\times8)$ deviations are resolved at all. The $-5069$~ppm of the plain-HMC
arm is the largest of them and is $1.1$ standard errors, so it is a fluctuation
rather than a defect of the map; the classical arms at the same volume are
resolved because their deviations are an order of magnitude larger.

\begin{table*}[htbp]
\centering
\small
\caption{U(2) preconditioner benchmark: relative deviation of each arm's mean
Wilson loop from the exact closed form, in parts per million, at both
benchmarked couplings, averaged over the settled half of $400$ trajectories on
$64$ chains. Parenthesised digits are the standard error on the preceding
deviation in the same units (\app{stats}); the hot-start rows are given to one significant
figure, their deviations exceeding their errors by two orders of magnitude or
more. The last column is $\langle Q^2\rangle$ as a ratio to the
exact $1.0012$. Arms are named by starting configuration and by which winding
update of \sect{instanton_updates_classical} is enabled.}
\label{tab:u2_observables}
\begin{tabular}{lrrrr@{\hspace{1.6em}}c}
\toprule
arm & $W(1{\times}1)$ & $W(2{\times}2)$ & $W(4{\times}4)$ & $W(8{\times}8)$
& $\langle Q^2\rangle/$exact \\
\midrule
\multicolumn{6}{l}{\textit{$L=32$, $\beta=105.651$}} \\
lift, plain HMC     & $+7(13)$ & $-9(63)$ & $-156(524)$ & $-5069(4676)$ & $1.14$ \\
cold, plain HMC     & $+71(12)$ & $+421(76)$ & $+3201(498)$ & $+2.5\times10^{4}$ & $0.00$ \\
hot, plain HMC      & $-5.5\times10^{4}$ & $-2.3\times10^{5}$ & $-6.5\times10^{5}$ & $-9.8\times10^{5}$ & $21.2$ \\
cold $+$ even-only & $+69(11)$ & $+197(61)$ & $+867(461)$ & $+4172(3917)$ & $0.82$ \\
lift $+$ even-only & $+18(12)$ & $+29(72)$ & $+422(546)$ & $+381(4346)$ & $1.05$ \\
hot $+$ even-only  & $-5.5\times10^{4}$ & $-2.3\times10^{5}$ & $-6.4\times10^{5}$ & $-9.4\times10^{5}$ & $1.10$ \\
cold $+$ marginal  & $+64(11)$ & $+149(65)$ & $+970(477)$ & $+6376(3501)$ & $1.01$ \\
lift $+$ marginal  & $+1(11)$ & $+1(71)$ & $-227(518)$ & $-1898(3547)$ & $0.98$ \\
\midrule
\multicolumn{6}{l}{\textit{$L=64$, $\beta=416.524$}} \\
lift, plain HMC     & $-1.1(2.2)$ & $+3.5(11.4)$ & $+0.5(82.8)$ & $+341(610)$ & $1.14$ \\
cold, plain HMC     & $+51(2)$ & $+229(10)$ & $+1218(75)$ & $+7644(571)$ & $0.00$ \\
hot, plain HMC      & $-6.3\times10^{4}$ & $-2.5\times10^{5}$ & $-7.0\times10^{5}$ & $-9.9\times10^{5}$ & $87.7$ \\
cold $+$ even-only & $+57(2)$ & $+247(10)$ & $+1273(65)$ & $+7115(479)$ & $0.89$ \\
lift $+$ even-only & $+0.2(1.9)$ & $+3.7(11.2)$ & $+129(80)$ & $+1263(510)$ & $1.03$ \\
hot $+$ even-only  & $-6.2\times10^{4}$ & $-2.4\times10^{5}$ & $-7.0\times10^{5}$ & $-9.9\times10^{5}$ & $1.16$ \\
cold $+$ marginal  & $+56(1)$ & $+233(8)$ & $+1182(64)$ & $+6873(476)$ & $0.98$ \\
lift $+$ marginal  & $-1.4(1.5)$ & $-10(9)$ & $-60(68)$ & $+395(498)$ & $1.00$ \\
\bottomrule
\end{tabular}
\end{table*}

A cold start under plain HMC and the same budget is $51$, $229$ and
$7644$~ppm high on those three loops. Comparing each cold arm with the
preconditioned arm that ran the same sampler, the cold start is $23$ to $68$
times further from exact on the plaquette and $W(2\times2)$; the one larger
ratio is $276$, under the even-only move, where the preconditioned plaquette
sits at $0.2$~ppm. Its topological
charge never moves in the $400$ trajectories, so its $\langle Q^2\rangle$ is
exactly zero times exact rather than merely small. Adding the winding move
repairs the charge and not the local observables: $\langle Q^2\rangle$ comes
back to $0.98\pm0.03$ of exact under the marginal move, while the plaquette
stays at $56$~ppm. A hot start does not begin to equilibrate within the
budget: it is $62{,}500$~ppm low on the plaquette and $994{,}000$~ppm low on
$W(8\times8)$. It carries $\langle Q^2\rangle$ at $88$ times exact.

\subsection{Distribution comparison}
\label{sec:distributions}

\begin{figure*}[t]
  \centering
  \includegraphics[width=\linewidth]{u2_2d/fig62_uv_distributions.png}
  \caption{Distributions of Wilson loops over individual configurations,
  compared with the exact result, for the preconditioned chain and for cold- and
  hot-started HMC with the same plain sampler: U(1) at $L=32$, $\beta=218.58$,
the hardest coupling of the U(1) comparison, and U(2) at $L=64$,
$\beta=416.524$, the per-arm benchmark of \tb{u2_observables}. Each configuration contributes one value, taken from the second
  half of every chain. The horizontal axis is the deviation from the exact value
  in units of $\sigma_{\rm cfg}$, the spread of the preconditioned distribution.
  Each panel gives $\sigma_{\rm cfg}$ and the width of each classical
  distribution relative to the preconditioned one. Distributions that fall
  outside the plotted range are shown as arrows labeled with their centre, in
  units of $\sigma_{\rm cfg}$.}
  \label{fig:uv_distributions}
\end{figure*}

\begin{figure*}[t]
  \centering
  \includegraphics[width=\linewidth]{u2_2d/fig63_topology_distributions.png}
  \caption{Distributions of the topological charge compared with the exact
  result, at a coupling where plain HMC cannot change the charge in either
  theory. The insets give the exact $\langle Q^2\rangle$ and each chain's
  $\langle Q^2\rangle$ as a ratio to it. The percentage is the fraction of the
  hot-start distribution that lies outside the plotted range.}
  \label{fig:topology_distributions}
\end{figure*}

The comparisons above are all of means, but a preconditioner should also approximately start with the correct observable distributions. \figs{uv_distributions}{topology_distributions} compare the per-configuration
distributions against the closed form, in both theories.

Distances here are quoted in units of $\sigma_{\rm cfg}$, the spread of a given
Wilson loop across individual configurations of the preconditioned ensemble (\app{stats}).
On loops from the plaquette out to $W(8\times8)$ the
preconditioned distribution is centered on exact to within
$0.05\,\sigma_{\rm cfg}$ at every loop size, and its width matches that of a
cold chain under the same sampler to within $3.2\%$, so the spread is the
theory's. For the cold start, the observables fluctuate correctly about an incorrect mean value. The hot start matches neither the mean value nor the distribution. At the largest loop it is
too narrow rather than too wide, that loop having collapsed to about a percent
of its exact value. On $P(Q)$ the preconditioned histogram tracks
the exact sector weights, while a cold start puts every configuration at $Q=0$
and a hot start scatters more than half its weight beyond the plotted range at
$L=64$.

\subsection{Cost accounting}
\label{sec:cost}
An interval measured on a chain that has not equilibrated counts decorrelation
within the topological sector it started in, not a supply of independent
configurations of the target, so $t_{\rm interval}$ is informative only on an arm
that has equilibrated. Local observables are unbiased inside that sector and
decorrelate quickly there,
which is why a plaquette interval looks short on precisely the arms that are
worst off; the ensemble average is nonetheless wrong, by the amount the sector
weight is wrong. Where a classical arm does not equilibrate within the budget we take its cost
as infinite and report the improvement factor as a lower bound, that bound being
the coupling's own trajectory budget (\tb{grid}) divided by the preconditioned
cost. What $F$ counts on both sides is HMC trajectories. The cost of building the
lifted configuration falls outside that count, and it is not small: in
U(1), where the classical baseline still converges, the campaign that produced
the checkpoint costs more than every classical burn-in that converges, and a
lifted configuration costs more than a trajectory, so measured in seconds the
method loses at every coupling where the classical baseline is correct; in U(2)
the top rung of the ladder costs roughly twice what a winding-HMC configuration
costs. $F$ is therefore not a wall-clock speed-up, and the case for the method rests on
the regime where no classical budget is sufficient rather than on cost in the
regime where one is (\app{overhead}).


\begin{figure*}[t]
  \centering
  \includegraphics[width=\linewidth]{u2_2d/fig64_money.png}
  \caption{Improvement factor over winding HMC, against
  coupling, at both evaluated volumes. Three checkpoints are drawn: the deployed
  \texttt{default}, \texttt{cov60}, which is trained densely below a low
  ceiling, and \texttt{wide\_dense}, which extends the ceiling to $\beta=2000$;
  \tb{volume_scaling} scores all four. Circles are measured values. Triangles are
  lower bounds, at couplings where no classical chain equilibrated within that
  coupling's own trajectory budget (\tb{grid}). Crosses mark couplings where the lift itself did not
  equilibrate and no factor exists; they carry no interval, a chain that never
  crosses having no crossing to resample. Dotted lines mark each checkpoint's
  training ceiling in its own colour; \texttt{wide\_dense}'s, at
  $\beta=2000$, lies beyond the right edge of the grid, so no test
  coupling is ever an extrapolation for it. Vertical bars are a chain-bootstrap interval, resampling the chains
  of each arm and propagating both timescales (\app{stats}); they are wide
  because $t_{\rm therm}$ is a first-crossing index set by the last chain to
  settle.}
  \label{fig:money}
\end{figure*}

\fig{money} measures the improvement factor against winding HMC across
the coupling range, at both evaluated volumes, and the result separates into
two regimes.

Every case in which winding HMC equilibrates lies below $\beta\approx90$, and
there the improvement factor is measured rather than bounded: eighteen such
cases, spanning six distinct couplings across the three checkpoints of
\fig{money} and the two volumes. It is modest, with a median of $1.4$, a range
of $0.1$ to $18$, and a value below one in six of the eighteen. Counted in
trajectories alone and against the strongest classical sampler this theory
admits, the preconditioner is break-even at best, which is the expected result
in a regime where a good classical sampler already exists. The converse does not
hold: winding HMC fails to equilibrate at two couplings below the crossover,
$\beta=41.5$ at $L=64$ and $\beta=58.0$ at $L=32$, which is why \fig{money}
carries lower bounds there as well.

Above $\beta\approx90$ no classical arm reaches equilibrium within the budget
at any coupling tested, so its cost is infinite by the gate above and the
improvement factor is a lower bound, running from $1.5$ to $200$ with a median
of $9.1$ over twenty-four cases. What fails there is not topology: the winding
move keeps firing, and even at $\beta=1660$ it holds $\langle Q^2\rangle$ within
$25\%$ of exact (\tb{u2_L128}). It is the local observables that never settle,
so the chain decorrelates without ever equilibrating.

Plain HMC is the weaker baseline and fails earlier. Its interval, maximized over
$Q^2$, is finite only at the two softest couplings of the grid, $\beta=10.9$ and
$13.9$ at $L=32$; from $\beta\approx25$ upward, and at every $L=64$ coupling
including the lowest, its charge does not move within the budget and it supplies
no independent configurations of the target at any cost.


\begin{figure*}[t]
  \centering
  \includegraphics[width=\linewidth]{u2_2d/fig65_ladder.png}
  \caption{Cost of one independent, correctly distributed configuration along
  the ladder. Each step doubles $L$ and approximately quadruples $\beta$, the
  fine coupling being fixed by \Eq{topology-matching}, which holds
  $\chi_t V$ constant, so the three rungs describe the same physical
  system at three lattice spacings. The first two rungs use the deployed
  checkpoint and the third uses the wider-coverage one, since only its training
  range reaches $\beta=1660$. (a) The number of trajectories the preconditioned
  chain needs to equilibrate, and the number the classical chains need between
  independent configurations; the preconditioned chain is already equilibrated
  at the first two rungs (dotted line). (b) The resulting improvement factor.
  Every value is a lower bound: arrows mark rungs where no classical chain
  equilibrated within the $400$-trajectory budget, and at the first rung the
  preconditioned cost of zero is counted as one trajectory. Vertical bars are
  the chain-bootstrap interval of \app{stats}, asymmetric on the preconditioned
  arm at the first two rungs for the reason given there.}
  \label{fig:ladder}
\end{figure*}

\subsection{Reach in volume}
\label{sec:multi-lift}

Reaching a target $(L,\beta)$ may take several lifts, and the error does not
compound when it does (\app{multilift}): it enters at the last rung, so
endpoint accuracy is set by that rung's distance from training coverage and not
by the number of rungs climbed. Laddering therefore extends the volume reached rather than the
coupling, since each lift multiplies $\beta$ by about four and the last one lands
at the same fine coupling whatever path reached it.

\begin{table*}[htbp]
\centering
\caption{U(2) preconditioner benchmark at $L=128$, $\beta=1660.08$: $32$ chains,
$400$ trajectories, against an exact plaquette of $0.99879505$ and an exact
$\langle Q^2\rangle$ of $1.0012$. Plaquette columns are relative deviations from
exact in parts per million, before any trajectory and at the last record;
$\langle Q^2\rangle$ is averaged over the last half of the history and quoted as
a ratio to exact. ``$P(Q)$ cov.'' is the exact probability of the sectors an arm
visits and ``odd'' the number of odd sectors it reaches.}
\label{tab:u2_L128}
\begin{tabular}{lrrr@{\,$\pm$\,}lrr}
\toprule
arm & plaq. $t=0$ & plaq. final & \multicolumn{2}{c}{$\langle Q^2\rangle/$exact}
& $P(Q)$ cov. & odd \\
\midrule
lift, plain HMC    & $-0.8$ & $-0.02$ & 1.155 & 0.309 & 0.941 & 3 \\
cold $+$ even-only   & $+1200$ & $+25$ & 0.786 & 0.049 & 0.507 & 0 \\
cold $+$ marginal    & $+1200$ & $+24$ & 1.249 & 0.052 & 1.000 & 4 \\
lift $+$ marginal    & $-0.8$ & $-0.14$ & 1.244 & 0.063 & 1.000 & 4 \\
\bottomrule
\end{tabular}
\end{table*}

The ladder was pushed one rung past the deployed configuration, to $L=128$ at
$\beta=1660.08$, four times the largest linear size in the training set and
sixteen times the sites (\fig{ladder}). Before a single trajectory the
lifted configuration sits $0.8$~ppm from the exact plaquette against a
cold start's $1200$~ppm, and it reaches equilibrium in $25$--$30$ trajectories.
Neither classical arm reaches equilibrium at all within the budget, so the
improvement factor there is a lower bound.

Two columns of \tb{u2_L128} have to be read together rather than separately. The
sector coverage is the exact probability of the sectors an arm visits, so an arm
that scatters charge over many sectors scores well on it whatever its variance;
only with $\langle Q^2\rangle$ beside it does coverage distinguish a correctly
populated ensemble from a badly spread one. And the charge interval on the arm
that runs plain HMC is wide because that sampler cannot change $Q$: its $32$
chains contribute $32$ frozen charges rather than a sampled history, so the
interval measures the spread of the lift's output and not the mixing of a chain.

The charge variance at this rung overshoots, and the overshoot is not a property
of the preconditioner. Both arms that run the marginal move land at
$\langle Q^2\rangle$ about a quarter above exact, $1.249\pm0.052$ from a cold
start and $1.244\pm0.063$ from a lift, four standard errors high and in
agreement with each other; the arm that starts from a lift and runs plain HMC,
which cannot change the charge it was given, sits at $1.155\pm0.309$ and is
consistent with exact. That the cold and the lifted marginal arms agree while
starting from opposite charge distributions points at the stationary
distribution of the move at this coupling rather than at the configurations
handed to it, and the arm whose charge really is purely transported is the one
that is consistent with exact. We have not isolated the cause.

\section{Training coverage}
\label{sec:coverage}

Where the preconditioner degrades is governed by coupling coverage: which
couplings a checkpoint was trained on, and how far the test coupling sits from
them. This section measures how far that reach extends, and what widening the
training range recovers. The same question has been asked of a conditional flow
for two-dimensional U(1), by interpolating its conditioning coupling to values
at which it was never trained and then re-training on the model's own samples to
reach further \cite{singha2023conditional}.

A checkpoint's \emph{training ceiling} is the largest coupling at which it was
ever given training data. Because the network is conditioned on the fine
side of each training pair, the ceiling is the largest coupling a checkpoint was
asked to generate \emph{at}, with the coarse side near $\beta_f/4$.

A checkpoint can be trained at a coupling where the classical sampler does not
deliver, because an ensemble is built once and can afford a burn-in no deployed
chain could (\app{settings}). That settles local structure but not topology, so
the charge content is manufactured by the augmentation mechanism of
\sect{naive-blocking}. Wrong sector frequencies cost nothing downstream, since the network is
asked only for local structure at a given charge and the charge itself is
transported at deployment.

The experiment in this section is retraining on different sets of couplings. We
call one of those a \emph{training rung}, by analogy with the ladder rungs of
\sect{method}, although a training set is a list of couplings and need not be a
ladder. \texttt{default}, \texttt{wide} and \texttt{wide\_dense} share
architecture, capacity, optimization and data recipe and differ only in which
fixed rungs their training sets contain, so any change among them is
attributable to coverage alone; \texttt{cov60} differs additionally as noted
below. Three U(1) checkpoints are compared: \texttt{deployed}, with a ceiling at
$\beta=60$ and the one used for every other U(1) result here; \texttt{wide2000},
which extends the ceiling to $\beta=2000$ with $11$ fixed rungs against
\texttt{deployed}'s $4$; and \texttt{wide2000\_dense}, which adds $30$ further
rungs inside that range without extending it. There are four U(2) checkpoints: \texttt{cov60} and \texttt{default}
are narrow, with ceilings at $\beta=227$ and $416$, \texttt{default} being the
one deployed for every other result here; \texttt{wide} extends the ceiling to
$\beta=2000$ using $23$ rungs against \texttt{default}'s $11$; and
\texttt{wide\_dense} adds $31$ further rungs inside that range without extending
it. \texttt{default}'s rungs are not spread over the coupling axis but placed on
the ladder: its training couplings are those the deployed ladder of
\sect{naive-blocking} visits.

Each is scored over a grid held fixed across checkpoints, listed in
\app{grid}, and run twice, once
with plain HMC in every arm and once with the marginal winding move of
\sect{instanton_updates_classical}, the stronger of the two, so preconditioned
and classical starts are always paired under an identical sampler. ``Winding
HMC'' refers to that move wherever U(2) results are reported.
\fig{money} draws the three that carry the two claims;
\tb{volume_scaling} scores all four and \tb{u1_coverage_tests} gives the
matching U(1) retrain, both in \app{tables}.

Range is the axis that moves the coupling at which the method stops working.
\texttt{default} stops equilibrating
above $\beta=415$ at $L=32$ and above $\beta=94$ at $L=64$, while
\texttt{wide\_dense} works at every coupling tested, to $\beta=1623$ and
$\beta=1660$ respectively. \texttt{wide} does not: at $L=64$ it fails at three couplings inside its own
trained range, $\beta=174$, $637$ and $1660$ under winding HMC and
$\beta=94$, $637$ and $1660$ under the plain one, so range alone does not
guarantee that the coupling reached moves with it.
Density acts below the ceiling instead, and much more weakly. Adding $31$ rungs inside
\texttt{wide}'s existing range improves it at $30$ of $44$ cases, but by a
median of $2.1\times$ where range gives one to two orders of magnitude. Pooled
over all $44$ cases that gain is resolved, $[1.2,2.7]$, while in either half
alone the interval spans unity, so the effect is small and is resolved only in the combined sample. In U(1) density changes nothing measurable, which we attribute to those training
sets already being dense: \texttt{wide2000} and \texttt{wide2000\_dense} carry
an identical block of $78$ randomly placed couplings below $\beta=250$ as well
as the same $11$ fixed rungs, so the $30$ rungs separating them are a small
increment on a grid already filled. That pair is a clean comparison; the U(1)
range comparison is not, because \texttt{wide2000} raises the ceiling of that
random block from $\beta=60$ to $250$ as well as adding fixed rungs, so two
coverage changes move together. Both widen coverage, so the direction is
unambiguous, but its size cannot be assigned to the fixed rungs alone. The U(2) checkpoints carry no random rungs, so there the same
comparison, of $31$ added rungs, is between a sparse grid and a dense one, and
the ambiguity does not arise.
\tb{u2_coverage_tests} carries both comparisons, with the caveats that qualify
their size. The two axes therefore act where the other does not, and on the
evidence here a fixed rung budget is better spent extending the range than
resampling what it already covers.

The curves in \fig{money} are irregular because $t_{\rm therm}$ is a discrete
first-crossing index, so the level of a curve carries more than any point on it.
Above
$\beta\approx90$ every point is also a lower bound, since no classical arm
equilibrates there and the plotted value is that coupling's budget divided by
the preconditioned cost alone: the peak at $L=64$, $\beta=174$ is one of these,
where the lift arrives already equilibrated against $t_{\rm therm}$ of $44$ and
$84$ trajectories at its neighbors. That there is no clean trend in $\beta$ is itself the
result: the factor tracks distance to the nearest training rung, and the rungs
sit at fixed, unevenly spaced couplings, so that distance falls and rises again
as $\beta$ passes each in turn.

Two limits bound what this grid can show. The first is geometric. Every grid
point is itself a matched lift, doubling $L$ and fixing $\beta_f$ by
\Eq{topology-matching}; what moves along the grid is not the lift but where it
lands, and because $L_f$ is held fixed while $\beta_f$ rises, successive points
target theories of smaller physical volume. The charge distribution narrows with
them: at $L=32$ the exact $\chi_t V$ falls from $0.22$ at $\beta=415$ to
$0.0008$ at $\beta=1623$, where $999$ configurations in $1000$ carry $Q=0$. At
the top of the grid there is almost no topology left to transport, and a cold
start, which carries $Q=0$ exactly, is already close to correct on that count;
what the improvement factor measures there is the equilibration of local
observables alone. The topological claim rests on the ladder, where
\Eq{topology-matching} holds $\chi_t V$ exactly constant and the physical volume
with it, and not on this grid.

That caps the grid at $\beta\approx1660$, below the wide checkpoints' ceiling of
$2000$, so no test coupling ever asks them to extrapolate: only the narrow pair
is tested outside its trained range, and there it fails at all seven couplings
above \texttt{default}'s ceiling. The comparison therefore establishes that a
checkpoint fails outside its own coverage, but not that widening helps for any
reason beyond bringing the test back inside it. Separating the two would need a
grid above $\beta=2000$, at $L\approx200$ to keep the physical volume from
collapsing, beyond what was run here.
\texttt{cov60} differs from the other three in three further settings: it was
trained without the random symmetry augmentation of \sect{scoring-network},
without the bias that makes the training sampler draw high noise levels more
often at large $\beta$, and without the $\beta$-aware floor on the noise
schedule. We measured whether those matter rather than assuming they do not. Retraining it with all three restored,
on the same rungs and the same budget, and scoring the raw lift at the same
couplings leaves it unchanged: the median ratio of their raw-lift deviations from exact is $1.00$, and the
retrained model is the closer of the two at $7$ of the $14$ couplings. The
contrast between \texttt{cov60} and \texttt{default} is therefore a coverage
effect. The independent check on the same axis is \texttt{wide\_dense} against
\texttt{wide}, which differ by nothing but the $31$ added rungs.

The second limit is that a ceiling predicts failure only roughly, and errs both
ways: \texttt{cov60} still equilibrates at $\beta=264$, above its own ceiling of
$227$, while \texttt{default} has already failed at $\beta=174$ and $L=64$,
under half of its $416$. The volume step moves that point without harming the
lift, whose deviation is smaller at $L=64$ for all four checkpoints
(\tb{volume_scaling}). The larger volume tightens the criterion rather than the map: an ensemble mean
over more sites carries a smaller error bar, so a bias of the same size takes
longer to fall within it.

\section{Outlook}
\label{sec:conclusions}
The objective of this program is four-dimensional SU(3) gauge theory and
ultimately full QCD with dynamical fermions, where the fermion determinant is
handled by a rational approximation to the Dirac kernel~\cite{kennedy1999} and
where topological freezing is a practical limit on the lattice spacings
reachable and where no exact reference
exists to score against. The two-dimensional theories studied here are a
controlled setting in which the method can be measured rather than the
destination.
The results above rest on two structural properties of the two-dimensional
theories: blocking preserves $Q$ exactly, and an inexpensive update exists that
changes it. Pure 4D SU(3) has neither, and differs in three
respects that bear on the method: its blocking map does not preserve $Q$, it has
no inexpensive topology-changing update, and its topological charge has no
unambiguous lattice definition. It is also not exactly solvable. That costs more than a scorecard: the
equilibration time of \sect{relaxation-correction} is defined as a crossing of
the bias against the exact value, so the closed form is the stopping rule as
well as the measure, and every ``equilibrated'' verdict in \figs{money}{ladder}
rests on it. In four dimensions a practitioner would have to replace it, by
agreement between arms started from different distributions, by a plateau
criterion, or by a long reference run in the role the closed form plays here.
The construction itself is unaffected, since nothing in the sampler consults the
exact value.

Transport needs two things, and the blocking map supplies only the first. It
pairs fine and coarse configurations while preserving $Q$ exactly, but only
because the plaquette phase is an Abelian sum of link variables, so the four
fine sub-plaquettes telescope into the coarse one; an ordered product of SU(3)
matrices does not telescope. The second requirement is distributional: the
coarse ensemble's $P(Q)$ must already be the fine theory's, since an exact map
applied to a charge from the wrong distribution reproduces the wrong distribution exactly. In two dimensions this holds automatically, because \Eq{topology-matching} holds $\chi_t V$ exactly
fixed from rung to rung and the deployed rungs inherit it
(\sect{naive-blocking}).
In 4D the principle survives at fixed physical volume, but neither the matching
nor the charge itself stays exact: the couplings would have to be matched
numerically by scale setting \cite{sommer1994,luscher2010}, and $Q$ has no cheap
unambiguous lattice definition. The field-theoretic form requires smoothing,
whose scheme and extent measurably affect the result at finite spacing
\cite{alexandrou2015}, while the index of a Ginsparg--Wilson operator is exact
but expensive \cite{hasenfratz1998}.
Imposing $Q_{\rm coarse}$ would then be exact only with respect to a chosen
definition.

The instanton field is the method's other structural tool, and it serves twice:
it manufactures sector coverage in the training data, and it imposes
$Q_{\rm coarse}$ on the lifted configuration. It exists in U(2) because the
group decomposes into a U(1) factor and a topologically trivial SU(2) factor,
with the charge carried entirely by the Abelian piece, so the U(1) construction
transfers directly. SU(3) does not have this property, and a cheap
topology-changing update for it is a central open question for this method.
Mechanisms that could supply one exist: parallel tempering on boundary conditions, introduced for the two-dimensional
$\mathrm{CP}^{N-1}$ model \cite{hasenbusch2017} and since carried to
four-dimensional $\mathrm{SU}(N)$ and to full QCD
\cite{bonanno2021,bonanno2024}; the related strategy of opening the boundary to
admit topological change and then restoring periodicity by a non-equilibrium
evolution, now driven by a stochastic normalizing flow, which has been scaled to
four-dimensional SU(3) \cite{bonanno2025su3,bonanno2026}; a stochastic dynamics able to jump
between sectors outright \cite{rothkopf2026}; or a deliberate bias of the
generating action, a chemical potential $-\mu Q$ or a bias potential in $Q$ as in
multicanonical sampling \cite{bergneuhaus1992} and metadynamics
\cite{laio2016,eichhorn2026}. We regard the last of these as the most promising,
since it is the cheapest and asks least of the mechanism. The present method asks less of any of them than a sampler would. They are needed only while
training data is built, and only to visit sectors: the model must learn the local
structure at a given charge, which a biased ensemble supplies correctly at each
$Q$ it reaches, while how often it reaches each one never enters a result,
because the charge is imposed by transport at deployment. Neither the acceptance
rate of such a mechanism nor the sector frequencies it produces need therefore be
correct, and no bias has to be reweighted away. What the closed form remains
necessary for is evaluation of the preconditioner.

The division of labor is the part of this method we expect to outlast the
particular architecture used here. Generating UV structure and
preserving IR and topological structure are different jobs, each best done by
the tool suited to it: the model supplies the short-distance structure and the
topological content, HMC certifies the former and cannot reach the latter. The same asymmetry has been
reported independently in two-dimensional $\phi^4$ theory, where a brief
rethermalization was found to remove a short-distance mismatch quickly and a
mismatch in a relevant direction slowly \cite{hasenfratz2026}, which suggests
that where the boundary between the two falls is a property of the blocking step
rather than of the theory or of the generative architecture. Normalizing flows
\cite{albergo2019,hasenfratz2026cascade}, gauge-equivariant L-CNNs
\cite{favoni2020} and other diffusion architectures
\cite{zhu2024,zhu2025,aarts2026,komijani2026} are natural candidates
for extending it to more complicated theories.

\section*{Data and code availability}
The complete implementation, the configuration files defining every checkpoint
and ladder reported here, and the analysis scripts that produce every figure and
table are available at
\url{https://github.com/sopanda0910/inverse_rg_testing}. The summary records
behind each table are included there; generated ensembles are regenerable from
the configurations and are not distributed.

Every computation reported here, comprising all training runs, ladders,
validation scans and coverage retrains, was performed on a single workstation
with an AMD Ryzen 7 260 (8 cores) and one NVIDIA RTX 5060 Laptop GPU with
$8$\,GiB.

%% FILL BEFORE CIRCULATING: funding, grant numbers and any thanks. Uncomment
%% the environment once it has content -- an empty acknowledgments block prints
%% a heading with nothing under it.
% \begin{acknowledgments}
% \end{acknowledgments}

\bibliographystyle{apsrev4-2}
\bibliography{main}

\clearpage

\appendix

\section{Network, training and sampler settings}
\label{app:settings}

The settings are as follows. The U(1) network has $324{,}954$ parameters at
hidden width $56$ and depth $4$; the U(2) one has $403{,}394$ at width $64$ and
the same depth; both use $3\times3$ kernels. Training is denoising score
matching under Adam, for $100$ epochs at batch size $16$ and learning rate
$3\times10^{-4}$ in U(1) and $120$ epochs at batch size $32$ and
$2\times10^{-4}$ in U(2), holding out a tenth of each training set for
validation; the deployed U(2) run stopped at epoch $109$ rather than $120$, and
both checkpoints hold the exponential moving average of the weights at their
best validation epoch. The U(1) objective carries one additional term, a soft
topological-charge penalty on the one-step denoised estimate at weight $0.1$;
the U(2) objective is denoising score matching alone. The
deployed U(2) checkpoint is trained on eleven fixed $(L,\beta)$ rungs spanning
nine distinct couplings, and the deployed U(1) one on four fixed rungs plus $78$
couplings drawn at random below its ceiling; \sect{coverage} varies exactly this list and
nothing else. A trajectory means one Metropolis-accepted integration of unit
length. The integrator step is scaled as $\beta^{-1/2}$ and the number of steps
inversely, holding that length fixed and the acceptance near $0.8$, so a
trajectory is not a fixed amount of work across the ladder: it is $9$ leapfrog
steps at $\beta=28$ and $72$ at $\beta=1660$. Improvement factors compare two
arms at one coupling and are unaffected, but trajectory counts should not be
compared between couplings without that factor.

Sampling uses $200$ ancestral steps, each followed by one
Langevin corrector step, with the charge projection applied every ten steps
below $\sigma=0.5$. The tail is $16$ local sweeps in U(1), and in U(2) $30$
conditional SU(2) heatbath sweeps followed by $10$ local sweeps, a local sweep
being one heatbath pass over the lattice followed by two overrelaxation passes.

The training ensembles themselves are burned in once per rung, at a cost no
deployed chain could afford: $200$ to $2000$ HMC trajectories in U(1), rising
with the coupling, and in U(2) sixty heatbath and microcanonical
overrelaxation \cite{creutz1987} sweeps followed by $150$ to $400$
trajectories.

\section{The analytic-score blend in U(2)}
\label{app:blend}

The analytic object the blend of \sect{reverse_diffusion} would contribute in
U(2) is the score of the determinant-sector marginal $w_{\det}$ of \Eq{wdet}. Probing it on the raw lift at $L=32$, at a
coarse coupling of $\beta_c=67.41$ and so a fine coupling of
$\beta_f=263.50$,\footnote{These probes are their own lifts rather than points
of the grid of \app{grid}, whose neighboring point runs from $\beta_c=67.59$ to
$\beta_f=264.2$.} the plaquette's relative deviation runs
$-6\times10^{-4}$, $-0.19$ and $-0.30$ at blend weights $0$, $0.5$ and $1$,
while the transported charge is destroyed at full weight: $\langle Q^2\rangle$
runs $0.45$, $0.47$ and $10.7$ against an exact $0.39$. The pattern repeats at
every coupling probed, including one lying on a training rung.

Why the U(1) construction does not carry over is not settled. One candidate is
that the noise-smeared coupling it uses is derived from the harmonic limit of a
Wilson weight, and $w_{\det}$ is not of Wilson form, so the smearing is applied
to a weight it was not constructed for. A second is that the lift needs the
conditional $p(\psi_f\mid\psi_c)$ while the analytic score is the $\psi$
marginal, and the two cannot be mixed at a fixed weight. We report the
measurement and leave the mechanism open.

\section{Exact observables in 2D U(1) and U(2)}
\label{app:exact}

\tb{u1-exact-observables} and \tb{u2-exact-observables} collect the closed forms
used for scoring, with the sources they follow.

\begin{table*}[htbp]
\centering
\small
\caption{Exact 2D U(1) observables (Wilson action). Closed forms after Migdal \cite{migdal1975} and Drouffe \& Zuber \cite{drouffezuber1983}; the topological charge distribution follows Bonati \& Rossi \cite{bonatirossi2019}; the finite-$V$ Wilson loop is cf.\ Rusakov \cite{rusakov1990} and Witten \cite{witten1991}.}
\label{tab:u1-exact-observables}
\begin{tabular}{p{4.5cm}p{11.5cm}}
\toprule
Observable & Closed form \\
\midrule
Character coefficient &
  $c_n(\beta) = \dfrac{1}{2\pi}\displaystyle\int_{-\pi}^{\pi} e^{\beta\cos\theta}\,e^{-in\theta}\,d\theta = I_n(\beta)$ \\
Normalized character coefficient &
  $r_n = c_n/c_0 = I_n(\beta)/I_0(\beta)$, so $r_1$ is the mean plaquette at infinite $V$ \\
Wilson loop, area $A$ (finite $V$) &
  $\langle W(A)\rangle = \dfrac{\sum_n r_n^{V-A}\,r_{n+1}^{A}}{\sum_n r_n^{V}}$ \\
Wilson loop (infinite $V$) &
  $\langle W(A)\rangle \to r_1^{A}$ \\
Topological charge distribution &
  $P(Q) \propto \displaystyle\int dk\; e^{-2\pi i k Q}\,\psi(k)^{V}$,\quad
  $\psi(k) = \dfrac{\int e^{\beta\cos\theta}\cos(k\theta)\,d\theta}{\int e^{\beta\cos\theta}\,d\theta}$ \\
Topological susceptibility &
  $\chi_t = \langle Q^2\rangle / V$ \\
\bottomrule
\end{tabular}
\end{table*}

\begin{table*}[htbp]
\centering
\small
\caption{Exact 2D U(2) observables. Irreps are labeled by $(j,k)$ with
$d_{(j,k)}=2j+1$, and $I_\nu$ is the modified Bessel function of order $\nu$.
After Migdal \cite{migdal1975}, with the partition function and finite-$V$
Wilson loop following Rusakov \cite{rusakov1990} and Witten \cite{witten1991},
the one-link integral Creutz \cite{creutz1980} and Kennedy \& Pendleton
\cite{kennedypendleton1985}, and $P(Q)$ Bonati \& Rossi \cite{bonatirossi2019}
applied to $w_{\rm det}$. The determinant-sector rows follow from the one-link
integral applied to the split representation; the minimum-KL matched coupling
is this work.}
\label{tab:u2-exact-observables}
\begin{tabular}{p{4.5cm}p{11.5cm}}
\toprule
Observable & Closed form \\
\midrule
Character coefficient &
  $c_{j,k}(\beta) = d_j \displaystyle\int_0^{2\pi}\dfrac{d\phi}{2\pi}\,e^{-ik\phi}\,g_j(\beta\cos\phi)$,\quad
  $g_j(a) = \dfrac{2\,I_{2j+1}(a)}{a}$ \\
Partition function &
  $Z = \displaystyle\sum_{(j,k)} \left(\dfrac{c_{j,k}}{d_j}\right)^{V}$,\quad
  the sum running over $j\in\tfrac12\mathbb{Z}_{\ge0}$ and $k\in\mathbb{Z}$ with
  $k\equiv 2j \pmod 2$ \\
Mean plaquette (infinite $V$) &
  $\langle\tfrac12\mathrm{Re\,Tr}\,P\rangle = \dfrac{I_1(x)\big(I_0(x)-I_2(x)\big)}{2\big(I_0(x)^2 - I_1(x)^2\big)}$,\ \ $x=\beta/2$ \\
Wilson loop, fundamental, area $A$ (infinite $V$) &
  $\langle\tfrac12\mathrm{Re\,Tr}\,W(A)\rangle = r_{\rm fund}^{A}$,\quad
  $r_{\rm fund} = \dfrac{c_{(1/2,1)}}{2\,c_{(0,0)}}$ \\
Wilson loop, fundamental (finite $V$) &
  $\langle\tfrac12\mathrm{Re\,Tr}\,W(A)\rangle = \dfrac{1}{2Z}\displaystyle\sum_{(j,k)}\ \sum_{\pm}\ \dfrac{d_j}{d_{j\pm1/2}}\left(\dfrac{c_{j,k}}{d_j}\right)^{\!A}\left(\dfrac{c_{j\pm1/2,\,k+1}}{d_{j\pm1/2}}\right)^{\!V-A}$,\ \ $d_j = 2j+1$ \\
& \footnotesize The fundamental fuses multiplicity-free, $(\tfrac12,1)\times(j,k) = (j{+}\tfrac12,k{+}1)\oplus(j{-}\tfrac12,k{+}1)$, so the double sum over $(r,r')$ collapses to two terms per irrep. Corrections to the area law are $O(e^{-\sigma(V-A)})$ from loops wrapping the torus, not a perimeter term \cite{rusakov1990,witten1991}. \\
\midrule
\multicolumn{2}{l}{\textit{Determinant sector} ($\psi=\mathrm{wrap}(2\theta)=\arg\det U$, an exact compact U(1) theory)} \\
\midrule
Single-plaquette marginal weight &
  $w_{\rm det}(\alpha) = \dfrac{2 I_1(z)}{z}$,\quad $z=\beta\cos(\alpha/2)$ \\
Topological charge distribution &
  $P(Q) \propto \displaystyle\int dk\, e^{-2\pi ikQ}\,\psi_{\rm det}(k)^V$,\quad
  $\psi_{\rm det}(k) = \dfrac{\int w_{\rm det}(\alpha)\cos(k\alpha)\,d\alpha}{\int w_{\rm det}(\alpha)\,d\alpha}$ \\
Topological susceptibility &
  $\chi_t = \langle Q^2\rangle/V$ (det.\ sector) \\
Matched U(1) coupling &
  $\beta_1(\beta)$: min.-KL projection s.t.\ $r_1(\beta_1) = \langle\cos\alpha\rangle_{w_{\rm det}}(\beta)$, with $r_1$ as above;\ \ $\to \beta/4$ as $\beta\to\infty$ \\
\bottomrule
\end{tabular}
\end{table*}

\section{Topology-matched ladder couplings}
\label{app:matching}

\Eq{topology-matching} has no closed-form solution, so $\beta_f$ is obtained
numerically: $\langle Q^2\rangle$ is evaluated from the character expansion of
\app{exact} at fixed $L$, and the root is bracketed between $\beta_c$ and
$64(\beta_c+1)$ and taken by Brent's method to a tolerance of $10^{-9}$. The
bracket is safe because $\langle Q^2\rangle$ at fixed $L$ decreases
monotonically in $\beta$, and the left end is below the solution for every step
of the form $L_f=2L_c$.

\tb{schedules} records the resulting schedule, and for reference the two
criteria it was chosen over, all three iterated from the same base ensemble. Two alternatives are shown with it. Matching the first character coefficient
$r_1$, which for this theory is the same as matching the mean plaquette, leaves
$\langle Q^2\rangle$ low by $2.2\%$ after one step and $2.8\%$ after two; the
tree-level rule $\beta_f=4\beta_c$ leaves it low by $5.8\%$ and $7.1\%$.
Matching \Eq{topology-matching} leaves it unchanged to the five digits shown.
In total variation between the top rung's exact $P(Q)$ and the base's, the three
give $1.2\times10^{-5}$, $5.7\times10^{-3}$ and $1.5\times10^{-2}$. The
criteria land on couplings $2$--$3\%$ apart, which is a different target theory
rather than an error in any of them; we choose the one that holds the quantity
the ladder transports.

\begin{table}[htbp]
\centering
\small
\caption{The three ways of choosing the fine coupling, each iterated twice from
the same base at $L=16$, $\beta=28$, where the exact
$\langle Q^2\rangle=1.00119$. The last column is the total variation distance
between that rung's exact $P(Q)$ and the base's, the distortion a transported
charge inherits. The first row is the deployed schedule. For this theory
matching $r_1$, the normalized first character coefficient, is the same as
matching the mean plaquette.}
\label{tab:schedules}
\begin{tabular}{llrrr}
\toprule
Criterion & Rung & $\beta_f$ & $\langle Q^2\rangle$ & $\|P_f-P_c\|_{\rm TV}$ \\
\midrule
$\langle Q^2\rangle$ matched & $L=32$ & $105.651$ & $1.00119$ & $1.2\times10^{-5}$ \\
                             & $L=64$ & $416.524$ & $1.00119$ & $1.2\times10^{-5}$ \\
\midrule
$r_1$ matched                & $L=32$ & $108.052$ & $0.97851$ & $4.6\times10^{-3}$ \\
                             & $L=64$ & $428.413$ & $0.97327$ & $5.7\times10^{-3}$ \\
\midrule
tree level, $4\beta_c$       & $L=32$ & $112.000$ & $0.94338$ & $1.2\times10^{-2}$ \\
                             & $L=64$ & $448.000$ & $0.93053$ & $1.5\times10^{-2}$ \\
\bottomrule
\end{tabular}
\end{table}

\section{The coupling grid}
\label{app:grid}

Every checkpoint in \sect{coverage} is scored on the same grid, and each point of
it is an independent lift: a coarse ensemble at $\beta_c$ is lifted to
$\beta_f$ at twice the linear size, and no charge is carried between points.
\tb{grid} lists them. The grid holds $L$ fixed and raises $\beta$, so it is not
a ladder and the invariance of \Eq{topology-matching} does not hold along it;
\sect{coverage} gives the consequence.

The budget shortens as the coupling rises, since the runs at the top of the grid
exist to establish that no classical arm equilibrates rather than to measure how
quickly one does. Each lower bound in \sect{cost} is quoted against the budget
of the coupling that produced it.

\begin{table}[htbp]
\centering
\small
\caption{The evaluation grid. $\beta_c$ is the coarse coupling each lift starts
from, $\beta_f$ the matched fine coupling it targets, and $n_{\rm traj}$ the
trajectory budget given to every arm at that coupling.}
\label{tab:grid}
\begin{tabular}{rrr@{\hspace{2.5em}}rrr}
\toprule
\multicolumn{3}{c}{$L_f=32$} & \multicolumn{3}{c}{$L_f=64$} \\
\cmidrule(lr){1-3}\cmidrule(lr){4-6}
$\beta_c$ & $\beta_f$ & $n_{\rm traj}$ & $\beta_c$ & $\beta_f$ & $n_{\rm traj}$ \\
\midrule
$4.05$ & $10.93$ & $400$ & $8.38$ & $25.29$ & $400$ \\
$5.28$ & $13.86$ & $400$ & $12.17$ & $41.46$ & $400$ \\
$8.22$ & $24.63$ & $400$ & $25.05$ & $93.79$ & $400$ \\
$12.78$ & $44.00$ & $400$ & $45.06$ & $174.0$ & $200$ \\
$16.19$ & $58.03$ & $400$ & $79.12$ & $310.4$ & $200$ \\
$23.80$ & $88.76$ & $400$ & $160.7$ & $636.8$ & $150$ \\
$33.46$ & $127.5$ & $200$ & $275.2$ & $1095$ & $150$ \\
$47.44$ & $183.6$ & $200$ & $416.5$ & $1660$ & $150$ \\
$67.59$ & $264.2$ & $200$ & & & \\
$105.2$ & $414.9$ & $200$ & & & \\
$135.9$ & $537.4$ & $200$ & & & \\
$199.2$ & $790.9$ & $150$ & & & \\
$267.9$ & $1065$ & $150$ & & & \\
$407.3$ & $1623$ & $150$ & & & \\
\bottomrule
\end{tabular}
\end{table}

\section{Statistical methodology}
\label{app:stats}

Integrated autocorrelation times are computed as specified in
\cite{madrassokal1988,wolff2004}, including Sokal's automatic windowing, and
$t_{\rm interval}=2\tau_{\rm int}$. The standard error on a mean over $N$
correlated measurements is then
$\sigma_{\rm mean}^2 = 2\tau_{\rm int}\,\sigma^2/N$, with $\sigma^2$ the
variance across measurements, and every deviation quoted in units of a standard
error uses that $\tau_{\rm int}$-aware form. The per-configuration spread
$\sigma_{\rm cfg}$ of \sect{distributions} is instead the plain standard
deviation of a Wilson loop across the individual configurations of an ensemble,
carrying no $\tau_{\rm int}$ factor because it is a width rather than an error
on a mean.

The equilibration time is read off the data as a crossing: $t_{\rm therm}$ is
the first trajectory after which $|\langle\mathcal{O}\rangle_t-\mathcal{O}_{\rm
exact}|$ stays within two standard errors for five consecutive measurements,
maximized over the Wilson loops. In the released analysis records this quantity
is stored as \texttt{t\_therm\_threshold\_old}; the field named
\texttt{t\_therm} there holds an exponential-fit estimator, which is not what is
reported here and which fails outright on the classical arms, whose approach to
equilibrium has no resolvable transient. Both quantities are quoted with a
chain-resampling interval, drawn by resampling the chains of an arm with
replacement and recomputing the quantity on each resample; we report the
$16$th and $84$th percentiles. For $t_{\rm interval}$ this is the usual
statement about a noisy estimator. For $t_{\rm therm}$ it carries more weight,
because a first-crossing index is set by the last chain to settle, so its
dominant uncertainty is which chains the ensemble happens to contain, and a
resample in which the crossing never occurs is recorded as such rather than
being dropped. The resulting interval is asymmetric wherever an arm is already
at the floor, equilibrating within the first record: resampling which chains
contribute can only delay that crossing and never bring it earlier, so the upper
edge reports how early a resampled set of chains fails the five-record rule
rather than giving a second estimate of the equilibration time.

The resampling unit is the chain. Every error bar on a sampled ensemble quantity
resamples whole chains. Beyond the usual underestimate for an autocorrelated
series \cite{wolff2004}, the case here fails outright: a topologically frozen
chain has zero within-chain variance in $Q$, so pooling over configurations
reports a vanishing error on an arm that has sampled nothing. A generated
ensemble is the exception, since a lift's configurations descend from
independent base chains and are themselves the independent unit; there the
configuration-level bootstrap and the naive standard error agree closely, and
the generated-ensemble error bars reported here are the bootstrap ones.

\section{Error accumulation over multiple lifts}
\label{app:multilift}

Because a target may be reached by more than one lift, the error a single lift
introduces must not accumulate along the chain of them, and it does not. Across
eight cases, comprising both theories, an endpoint inside and one beyond the
training range, and intermediate local sweeps between rungs enabled and
disabled, three lifts to a fixed endpoint leave $0.84$--$1.02$ times the
deviation from exact that a single lift to the same endpoint leaves. The error
is instead introduced by the final rung, which is why endpoint accuracy tracks that rung's
distance from training coverage rather than the length of the ladder below it.

\section{Cost of building a lifted configuration}
\label{app:overhead}

The cost that \Eq{improvement-factor} does not count is producing the
lifted configuration: a base ensemble at the bottom of the ladder, plus
one reverse diffusion per lift. \sect{cost} gives its size, which on this
hardware is not small. It is dominated by the reverse-diffusion sampler, whose
step count is set for accuracy rather than speed and is the obvious place to
look for a reduction.

Two features of that overhead are what keep the trajectory comparison
meaningful. The base ensemble is generated once, at a coupling where sampling is
cheap, and is amortized over every configuration lifted from it; and the reverse
diffusion is a fixed number of network evaluations, the same at every coupling,
while $t_{\rm interval}$ in the numerator grows without bound as the classical
chain freezes. A fixed overhead cannot close a diverging gap, and above
$\beta\approx90$, where no classical arm equilibrates at all, there is no finite
quantity for it to close.

\section{Supporting tables}
\label{app:tables}

This appendix collects the coverage comparisons of \sect{coverage} and the
classical entry cost of \sect{u1obs}: the U(1) entry cost (\tb{u1_h2h}) and coverage
retrain (\tb{u1_coverage_tests}), then the U(2) coverage comparisons
(\tbs{u2_coverage_tests}{volume_scaling}).

\begin{table*}[htbp]
\centering
\caption{U(1): HMC with the winding update, $L=32$, cold-started except at
$\beta_f=4.44$, which is hot-started. Entry cost is
the burn-in, in trajectories, at which the quality verdict is read; quality is
whether every Wilson observable sits within $2.5$ standard errors of exact,
judged separately from $\langle Q^2\rangle$ because the winding update samples
topology correctly well before local observables equilibrate. The window is
looser than the two standard errors of $t_{\rm therm}$ (\app{stats}) because it
is a single pass-or-fail verdict on a whole set of observables at one budget
rather than a crossing index on one of them, so a stricter window would fail on
the worst of the set by chance. The
$\langle Q^2\rangle$ column is the $500$-trajectory value on every row,
including the two whose entry cost is read at $8000$. The preconditioned chain's
entry cost is zero at all four couplings.}
\label{tab:u1_h2h}
\begin{tabular}{lrrrl}
\toprule
$\beta_f$ & entry cost (traj) & $\langle Q^2\rangle$ & exact $\langle Q^2\rangle$ & quality \\
\midrule
4.44   & 500          & $7.08 \pm 0.38$  & 6.79  & pass \\
14.15  & 500          & $1.87 \pm 0.058$ & 1.90  & pass \\
55.02  & 500 / 2000 / 8000 & $0.47 \pm 0.008$ & 0.474 & fail at 500 and 2000, pass at 8000 \\
218.58 & 500 / 2000 / 8000 & $0.032 \pm 0.002$ & 0.029 & fail at all three \\
\bottomrule
\end{tabular}
\end{table*}

\begin{table*}[htbp]
\centering
\small
\caption{U(1) coverage comparisons on the matched $L{=}8\to16$ grid. Endpoint
and columns are as in \tb{u2_coverage_tests}. All fifteen couplings are past
\texttt{deployed}'s training ceiling; ``off-rung'' is the pre-registered
subgroup of eight placed in the gaps between \texttt{wide2000}'s rungs.
Widening the range gives a factor of $37$--$50$ here on the parts-per-million
endpoint, or $8$--$9$ on the standard-error endpoint of
\tb{u2_coverage_tests}; either way it is the U(1) counterpart of the U(2) split
in \sect{coverage}.}
\label{tab:u1_coverage_tests}
\begin{tabular}{llrrl}
\toprule
comparison & subset & $n$ & wins & median ratio \\
\midrule
\texttt{wide2000} vs \texttt{deployed} & all & 15 & 15 & $45.3\times$ \\
\texttt{wide2000\_dense} vs \texttt{deployed} & all & 15 & 15 & $50.4\times$ \\
\texttt{wide2000} vs \texttt{deployed} & off-rung & 8 & 8 & $37.4\times$ \\
\texttt{wide2000\_dense} vs \texttt{deployed} & off-rung & 8 & 8 & $45.2\times$ \\
\addlinespace
\texttt{wide2000\_dense} vs \texttt{wide2000} & all & 15 & 9 & $1.1\times$ \\
\texttt{wide2000\_dense} vs \texttt{wide2000} & off-rung & 8 & 4 & $1.1\times$ \\
\bottomrule
\end{tabular}
\end{table*}

\begin{table*}[htbp]
\centering
\small
\caption{U(2) coverage comparisons, paired over identical evaluation cases on a
grid of $14$ couplings at $L=32$ and $8$ at $L=64$, each run twice. The endpoint
is the raw lift's deviation from exact before any trajectory, worst case over
the plaquette, $W(2\times2)$ and $W(4\times4)$, in parts per million. ``Wins''
counts the couplings on which the first-named checkpoint is the closer of the
two; ``median ratio'' is the median of (second)/(first), so $R$ reads as ``the
first is $R$ times closer'', with a percentile bootstrap interval
\cite{efron1993}. Two qualifications on the range-against-density comparison of
\sect{coverage}: its two medians cover different subsets, so they bound the
separation rather than measure it; and on the endpoint pre-registered for the
U(1) arm, the worst-case deviation in units of its own standard error, the
same two comparisons give $9.8\times$ and $17.1\times$ rather than $34.7\times$
and $134.5\times$, with every ordering unchanged.}
\label{tab:u2_coverage_tests}
\begin{tabular}{llrrll}
\toprule
comparison & subset & $n$ & wins & median ratio & 95\% CI \\
\midrule
\texttt{wide} vs \texttt{default} & past ceiling & 14 & 14 & $34.7\times$ & $[19.4,93.3]$ \\
\texttt{wide\_dense} vs \texttt{default} & past ceiling & 14 & 14 & $134.5\times$ & $[21.9,203.5]$ \\
\texttt{wide\_dense} vs \texttt{wide} & past ceiling & 14 & 10 & $2.2\times$ & $[0.6,4.1]$ \\
\addlinespace
\texttt{wide} vs \texttt{default} & in coverage & 30 & 20 & $1.4\times$ & $[0.99,3.6]$ \\
\texttt{cov60} vs \texttt{default} & in coverage & 30 & 21 & $2.4\times$ & $[1.6,4.7]$ \\
\texttt{wide\_dense} vs \texttt{wide} & in coverage & 30 & 20 & $2.1\times$ & $[0.95,2.6]$ \\
\texttt{wide\_dense} vs \texttt{wide} & all & 44 & 30 & $2.1\times$ & $[1.2,2.7]$ \\
\bottomrule
\end{tabular}
\end{table*}

\begin{table*}[htbp]
\centering
\small
\caption{U(2): raw preconditioner quality by checkpoint and volume, over the
same grid ($14$ couplings $\times\,2$ rounds at $L=32$, $8\times2$ at $L=64$).
Entries are the median over couplings of the endpoint of
\tb{u2_coverage_tests}, in parts per million; lower is better, and no
checkpoint degrades with volume.}
\label{tab:volume_scaling}
\begin{tabular}{lcc}
\toprule
 & \multicolumn{2}{c}{raw lift deviation (ppm)} \\
\cmidrule(lr){2-3}
checkpoint & $L=32$ & $L=64$ \\
\midrule
\texttt{cov60}   & 8315 & 7115 \\
\texttt{default} & 8851 & 7422 \\
\texttt{wide}    & 2148 & 1812 \\
\texttt{wide\_dense} & \textbf{1272} & \textbf{702} \\
\bottomrule
\end{tabular}
\end{table*}

\end{document}